\documentclass[journal]{IEEEtran}
\usepackage{amsmath,amssymb,mathtools}
\usepackage{amsthm}
\usepackage{booktabs,array}
\usepackage{enumitem}
\usepackage{algorithm}
\usepackage{algpseudocode}
\usepackage{float}
\usepackage{microtype}
\usepackage[hidelinks]{hyperref}
\newtheoremstyle{plainn}{}{}{\itshape}{}{\bfseries}{.}{.5em}%
  {\thmname{#1}\thmnumber{ #2}\thmnote{ (#3)}}
\theoremstyle{plainn}

\newtheorem{thm}{Theorem}
\newtheorem{prop}[thm]{Proposition}
\newtheorem{lem}[thm]{Lemma}
\newtheorem{cor}[thm]{Corollary}
\newtheorem{claim}[thm]{Claim}
\newtheorem{defn}{Definition}
\theoremstyle{remark}
\newtheorem{remark}{Remark}
\renewcommand{\qedsymbol}{QED}
\newcommand{\row}{\operatorname{row}}
\newcommand{\rank}{\operatorname{rank}}
\newcommand{\Ann}{\operatorname{Ann}}
\newcommand{\Gr}{\operatorname{Gr}}
\newcommand{\Prob}{\operatorname{Prob}}
\newcommand{\poly}{\operatorname{poly}}
\setlength{\emergencystretch}{1.5em}

\begin{document}

\title{Image-Space Reduction and Polynomial-Time Feasibility for Functional Observers
with a Two-Dimensional Compressed Kernel}
\author{}
\maketitle

\begin{abstract}
For a linear system $(A,C)$ and a prescribed observer row block $L_0$, feasibility of a functional
observer of order $q$ is governed by two rank conditions on the pencil $\lambda H-H_A$, where
$H=\mathcal L N$, $H_A=\mathcal L AN$ and $N$ spans the unobservable subspace $K_0=\ker[C;CA]$. The layer index of a design
is $c=k-\rank H$, $k=\dim K_0$; the case $c=1$ was settled in earlier work, where the admissible kernel
directions are the real eigenvectors of the compression $M=\pi_{K_0}A|_{K_0}$.

This paper treats the case $c=2$ for the split problem. Our main result is a \emph{reduction of the fibre
variable}: for a fixed kernel direction $D$ the pencil depends on the fibre $W$ only through its image
$\Phi(W)=\varphi(W)\subseteq\varphi(T_D)\subseteq D^\circ$; the pencil's whole row space lies in
$D^\circ$, of dimension $\omega$, but $\Phi(W)$ itself lies in the image $\varphi(T_D)$, whose dimension
is at most $n-k$ (at most twice the number of output rows). Consequently the search over a Grassmannian
of arbitrary dimension collapses to a search over subspaces of this low-dimensional image space. We give
the exact realisability conditions for a prescribed $\Phi$, the resulting reduced feasibility criterion,
and a classification of the invariant-plane base under simple spectrum over $\mathbb C$ --- with the real
and the conjugate-complex branches treated separately --- which is the finite-enumeration tool that
makes the reduction effective. This yields an exact yes/no decision algorithm: for the split problem,
under simple spectrum over $\mathbb C$ and a fixed output dimension, it runs in bit complexity polynomial
in $n$ and in the bit length $\tau$ of the whole input $(A,C,L_0)$. We also show that the layer index is
\emph{not monotone} at fixed order: there are instances where the $c=2$ layer is feasible while the
$c=1$ layer is not. All structural claims are stated for the split problem; the unrestricted problem,
and the split problem at $c=2$ without simple spectrum over $\mathbb C$ or without a fixed output
dimension, are left open.
\end{abstract}

\begin{IEEEkeywords}
Functional observers, minimum order, compressed pencil, invariant subspaces, Grassmann bundles,
subresultants, bit complexity.
\end{IEEEkeywords}

\section{Introduction}\label{sec:intro}

This section fixes the problem, records what the companion papers settle and what they leave open,
lists the contributions, and closes with the notation used throughout.

\emph{Setting.} The pair $(A,C)$ is given, together with a prescribed row block $L_0$. A design of order
$q$ is a row block $\mathcal L=[L_0;R]$ with $q$ rows; its row space $V$ has dimension $q$. Write
$K_0=\ker[C;CA]$, $k=\dim K_0$, fix a basis $N$ of $K_0$, and put $H=\mathcal L N$, $H_A=\mathcal L AN$. The design is
\emph{feasible} when
\begin{equation}\label{eq:5a}
\row H_A\subseteq\row H,
\end{equation}
\begin{equation}\label{eq:5b}
\rank(\lambda H-H_A)=\omega\quad\text{for all }\lambda\in\mathbb C,
\end{equation}
with $\omega=\rank H$. Condition~\eqref{eq:5b} is equivalent to the vanishing of the gcd of all
$\omega\times\omega$ maximal minors of the pencil (a standard polarity/resultant fact; see
Section~\ref{sec:admissible} and \cite{vzgg13}).

\emph{Reading the title.} The phrase \emph{two-dimensional compressed kernel} in the title names the layer
index and nothing else: it means that the compressed pencil $\lambda H-H_A$ has kernel deficiency
$c=k-\omega=2$, i.e.\ $\rank H=\omega=k-2$ and $\dim\ker H=c=2$. It does not mean that the compression
$M=\pi_{K_0}A|_{K_0}$ has a two-dimensional kernel: $M$ may be invertible, and $\dim\ker M$ is unrelated
to $c$. The kernel direction $D=\ker H$ over which the paper quantifies is the $c$-plane of the layer,
and $c=2$ is the case in which that direction is a plane.

\emph{Prior work and what is new.} The existence question, the order spectrum and the minimum-order
characterisation were settled in \cite{companion1}, \cite{companion2} (see also \cite{fer26a},
\cite{fer26b}): for every order $q$ there is an exact criterion, and the feasible orders are monotone and
saturate. The \emph{complexity} of the general decision problem remains open. In \cite{companion3} the
case $c=1$ was solved structurally: the admissible kernel directions are the real eigenvectors of the
compression $M=\pi_{K_0}A|_{K_0}$, restricted by a prescribed-row orthogonality condition; the design
family is a finite union of relative Grassmann bundles; and one function-field gcd per chart decides the
whole chart.

\emph{Contributions.} This paper treats $c=2$. The headline is the image-space reduction stated below as
the third contribution: it is what collapses the search, and it is proved for every $c$. The
invariant-plane classification is the finite-enumeration tool that makes the reduction effective at
$c=2$. The contributions are:
\begin{enumerate}[label=\arabic*.,leftmargin=*,itemsep=1pt,topsep=2pt]
\item \emph{Layer bookkeeping (Section~\ref{sec:layers}).} The layer index is $c=k-\omega$ and the
``ambient row excess'' is $t=q-\omega$; the identity $c-t=k-q$ shows that $c$ and $t$ move together when
$q$ is fixed. A previous, ambiguous identification of $\omega$ with $\dim(V\cap K_0)$ is corrected:
$\omega=\rank H=k-\dim(K_0\cap V^\perp)$, while $\dim(V\cap K_0)=\omega$ holds only under the split
condition. We
separate the unrestricted problem \textbf{U} from the $K_0$-split problem \textbf{S} and show by a
counterexample that the split hypothesis cannot be derived from \eqref{eq:5a}.
\item \emph{The two fixed-kernel fibres (Section~\ref{sec:fibres}).} For fixed $D=\ker H$ we give the
exact fibre in both problems, including the exactness condition $K_0\cap(D+AD)=D$ in \textbf{U}, and the
\emph{two independent} prescribed-row conditions in \textbf{S} --- one on the base
($P_{K_0}(V_0)\subseteq D^{\perp_{K_0}}$) and one on the fibre ($W_0\subseteq W$).
\item \emph{The headline: image-space reduction (Section~\ref{sec:reduction}).} The fibre acts only
through the image $\Phi(W)=\varphi(W)\subseteq\varphi(T_D)\subseteq D^\circ$: the whole pencil lies in
$D^\circ$, of dimension $\omega$, while $\varphi(T_D)$ has dimension at most $n-k\le 2p$, so the search
over a Grassmannian of arbitrary dimension collapses to a search over subspaces of this low-dimensional
image space. The reduction uses no property of $c$; $c=2$ enters only when the base is enumerated. We give
the general proof, the exact realisability inequalities for a prescribed $\Phi$, and the reduced
feasibility criterion.
\item \emph{The finite-enumeration tool at $c=2$ (Sections~\ref{sec:base}--\ref{sec:decision}).} The
reduction turns the fibre into a subspace of $\varphi(T_D)$; making it effective needs the base to be
finite, and that is what simple spectrum over $\mathbb C$ supplies. The set $\mathcal I_2(M)$ is then
finite, of $O(k^2)$ points, split into a real-real and a conjugate-complex branch, and it is enumerated
explicitly. The image space has at most one projective-line component, and on a chart every real point is
admissible, so the generic-gcd criterion is sharp in both directions and the algorithm is an exact yes/no
decision procedure.
\item \emph{Bit complexity (Section~\ref{sec:complexity}).} For the split problem, under simple spectrum
over $\mathbb C$ and with the output dimension $p$ held fixed, the exact yes/no decision runs in
$\poly(n,\tau)$ bit operations, with arithmetic in number fields of degree at most $k^2$. Here $\tau$ is
the bit length of the whole input: the entries of $A$ and $C$ \emph{and} those of the prescribed block
$L_0$, without which $V_0$, $W_0$, $\Phi_0$ and $\mathcal B_{\rm pres}$ are not even determined.
\item \emph{A non-monotonicity example (Section~\ref{sec:layers:c}).} At one and the same order, the
$c=2$ layer can be feasible while the $c=1$ layer is not. This rules out ``only small layers matter'' as
a complexity strategy.
\end{enumerate}

\emph{Approach and structure.} The approach is uniform, and its organising principle is the image-space
reduction: a statement about the fibre $W$ is first pushed to its image $\Phi(W)$, and the resulting
condition is then read on finitely many charts, one per subspace of $D^\circ$. Section~\ref{sec:layers}
fixes the layer parameters and the split hypothesis.
Section~\ref{sec:fibres} computes the fixed-kernel fibres in the two problems.
Section~\ref{sec:reduction} proves the image-space reduction and the reduced criterion.
Section~\ref{sec:base} classifies the base at $c=2$ under simple spectrum over $\mathbb C$ and separates
the structural layer that is valid for every $c$.
Section~\ref{sec:decision} gives the algorithm, its correctness and its bit cost, and closes with a
worked example (Section~\ref{sec:example}).
Section~\ref{sec:conclusion} consolidates the limitations and the evidence.
Appendix~\ref{app:height} holds the height-propagation chain used in the complexity bound.

\emph{Notation.} $E=\mathbb R^n$ throughout, with the Euclidean inner product. $\pi_{K_0}$ is the
orthogonal projection, $K_0^\perp$ its kernel, and $K_0^*$ the dual space of $K_0$. For a subspace
$J\le K_0$ we write
\[
J^\circ:=\{\varphi\in K_0^*:\ \varphi|_J=0\}\le K_0^*
\]
for its annihilator in $K_0^*$, so that $\dim J^\circ=k-\dim J$; the case used
throughout is $J=D=\ker H$, for which $\dim D=c=k-\omega$ and hence $\dim D^\circ=k-c=\omega$. $\Gr(d,X)$ is the
Grassmannian of $d$-planes in $X$. All matrices are real; $p$ is the number of output rows and $q$ the
order of the design, while $L_0$ is the prescribed block and $V_0=\row L_0$.

\emph{Input length.} The input consists of the matrices $A$, $C$ and $L_0$, together with the order $q$ and the output
dimension $p$. We write $\tau$ for the maximum bit length of an entry of $A$, of $C$ or of $L_0$, so that
the input is encoded by $O(n^2)$ rational numbers of bit length at most $\tau$, i.e.\ by $O(n^2\tau)$ bits,
and every complexity statement in this paper is a bound in $n$ and $\tau$ alone. The prescribed block is
covered by $\tau$ on purpose: $V_0=\row L_0$ and the derived data $W_0=P_{K_0^\perp}(V_0)$,
$w_0=\dim W_0$, $\Phi_0=\varphi(W_0)$, $t=q-\omega$ and the set $\mathcal B_{\rm pres}$ of \eqref{eq:bpres}
all depend on $L_0$, so a bound in the bit length of $A$ and $C$ alone would not be a bound on the running
time of the algorithm. For a subspace $X$ we write $\dim X$ for its dimension.

\section{The layer parameters}\label{sec:layers}

This section fixes the three integer parameters attached to a design, isolates the split hypothesis under
which the structural results are stated, and records a counterexample showing that the layer index is not
monotone in the order. Parts A and B are bookkeeping; Part C is the reason the later sections treat $c$
as given rather than assuming the smallest layer.

\subsection{Three quantities that must not be confused}\label{sec:layers:a}

For a design $\mathcal L$ with row space $V$ of dimension $q$, define
\begin{equation}\label{eq:params}
\omega:=\rank H,\qquad c:=k-\omega,\qquad t:=q-\omega .
\end{equation}

\begin{claim}[layer identities]\label{clm:identities}
Let $D=\ker H=K_0\cap V^\perp$. Then
\begin{enumerate}[label=(\arabic*),leftmargin=*,itemsep=1pt,topsep=2pt]
\item {\raggedright $\omega=\rank H=\rank(\mathcal L N)=\dim\pi_{K_0}(V)
=k-\dim(K_0\cap V^\perp)=k-\dim D$;\par}
\item $c=k-\omega=\dim D$;
\item $t=q-\omega=\dim(V\cap K_0^\perp)$, and $q=\omega+t$;
\item in the split case $\dim(V\cap K_0)=\omega$, and then
$V=(V\cap K_0)\oplus(V\cap K_0^\perp)$; in general $\dim(V\cap K_0)$ is a third quantity, not determined
by $\omega$ and $t$ (Claim~\ref{clm:notnormal} below has $\dim(V\cap K_0)=1$ with $c=2$ and $t=3$);
\item $\dim D^\circ=k-\dim D=k-c=\omega$.
\end{enumerate}
Moreover $c-t=k-q$.
\end{claim}

\begin{proof}
For (1): $\ker(\mathcal L N)=\{u\in K_0:\mathcal L u=0\}$, i.e.\
$K_0\cap\ker\mathcal L=K_0\cap V^\perp$ (when $\mathcal L$ has full row rank, $\ker\mathcal L=V^\perp$);
hence $\omega=\rank H=k-\dim(K_0\cap V^\perp)$. On the other hand, the restriction map
$r:V\to K_0^*$, $v\mapsto\langle v,\cdot\rangle|_{K_0}$, has kernel $V\cap K_0^\perp$, so
$\rank r=q-\dim(V\cap K_0^\perp)=\dim\pi_{K_0}(V)$. Its adjoint $r^*:K_0\to V^*$,
$x\mapsto\langle x,\cdot\rangle|_V$, is the transpose of $\pi_{K_0}|_V$ and has kernel $K_0\cap V^\perp$;
hence $\dim\pi_{K_0}(V)=\rank r=\rank r^*=k-\dim(K_0\cap V^\perp)$. All these
quantities coincide. Part (2) is the definition of $c$ together with (1). Part (3) is the first
computation above, and $q=\omega+t$ is the definition of $t$. Part (4) is the definition of the split
condition (Definition~\ref{defn:split}); the verification of the counterexample's numbers is part of the
proof of Claim~\ref{clm:notnormal}. Part (5) is $\dim D^\circ=k-\dim D=k-c=\omega$, and
$c-t=(k-\omega)-(q-\omega)=k-q$.
\end{proof}

\begin{remark}\label{rem:hyp}
In the hyperplane layer $q=k$ one has $c=t=k-\omega$, so $c$ and $t$ coincide there; this is why the
distinction between them was invisible in \cite{companion3}. It does not make $\omega$ coincide with
either: $\omega=k-c$, and $\dim(V\cap K_0)=\omega$ holds only in the split case. For $q\ne k$ even $c$
and $t$ differ, by $c-t=k-q$.
\end{remark}

Thus $c$ and $t$ coincide only in the hyperplane layer $q=k$; outside it $c\ne t$ by $c-t=k-q$. The layer
is indexed by $c=k-\omega$, not by $\dim(V\cap K_0)$, and $\dim D^\circ=\omega$ is a different quantity
from both.

\subsection{The split hypothesis}\label{sec:layers:b}

\begin{defn}\label{defn:split}
A design is \textbf{$K_0$-split} (problem \textbf{S}) when $\dim(V\cap K_0)=\omega$; equivalently
$V=(V\cap K_0)\oplus(V\cap K_0^\perp)$. Designs satisfying \eqref{eq:5a}, \eqref{eq:5b} without this
condition form the \textbf{unrestricted} problem \textbf{U}.
\end{defn}

For $c\ge0$ put $\mathcal I_c(M):=\{D\le K_0:\ \dim D=c,\ M(D)\subseteq D\}$; in the case $c=2$ we write
$\mathcal I_2(M)=\mathcal I_c(M)$.

\begin{claim}[invariant kernel under split]\label{clm:invariant}
Let $\mathcal L$ be split and satisfy \eqref{eq:5a}, and put $D=\ker H=K_0\cap V^\perp$,
$M=\pi_{K_0}A|_{K_0}$. Then $M(D)\subseteq D$.
\end{claim}

\begin{proof}
By \eqref{eq:5a}, $\ker H\subseteq\ker H_A$, i.e.\ $AD\perp V$. Let $\varphi_d(u)=\langle u,d\rangle$ and
$\varphi_{Ad}(u)=\langle u,Ad\rangle$ for $u\in K_0$. Both vanish on $V\cap K_0$: the first because
$D\subseteq V^\perp$, the second by $V\perp AD$. Since the design is split, the annihilator
$(V\cap K_0)^\circ\subseteq K_0^*$ has dimension $k-\dim(V\cap K_0)=k-\omega=c$, and the $\varphi_d$,
$d\in D$, are $c$ independent elements of it; hence $\varphi_{Ad}\in\operatorname{span}\{\varphi_d\}$ and
$\pi_{K_0}(Ad)\in D$.
\end{proof}

\begin{claim}[the split hypothesis is not a normalisation]\label{clm:notnormal}
There is a design satisfying \eqref{eq:5a}, with exact kernel and $c=2$, whose $D=\ker H$ is not
$M$-invariant; the design is then not $K_0$-split.
\end{claim}

\begin{proof}[Construction]
Let $e_1,\dots,e_8$ be the standard basis of $E=\mathbb R^8$ and put
\[
K_0:=\operatorname{span}\{e_5,\dots,e_8\},\qquad f_i:=e_{4+i}\quad(i=1,\dots,4),
\]
so that $f_1,\dots,f_4$ is a basis of $K_0$ and $e_1,\dots,e_4$ is a basis of $K_0^\perp$. All vectors of
$K_0$ below are written in the $f$-basis. Let
\[
D:=\operatorname{span}\{f_1+f_2,\ f_3+f_4\}\le K_0,\qquad \dim D=2 .
\]
Fix the compression $M=\operatorname{diag}(1,1,2,3)$ in the $f$-basis and the coupling $A_{21}$ below,
whose rows are indexed by $e_1,\dots,e_4$ and whose columns by $f_1,\dots,f_4$:
\[
A_{21}:=\begin{pmatrix}
1&-1&0&0\\
0&0&1&0\\
0&0&0&0\\
0&0&0&0
\end{pmatrix},
\]
and define $A$ on the basis $(f_1,\dots,f_4,e_1,\dots,e_4)$ of $E$ by
\begin{equation}\label{eq:notnormal-A}
\begin{gathered}
A=\begin{pmatrix} A_{11} & A_{12}\\ A_{21} & A_{22}\end{pmatrix}
=\begin{pmatrix} M & 0\\ A_{21} & 0\end{pmatrix},\\
Af_i=Mf_i+A_{21}f_i,\qquad Ae_j=0\quad(j=1,\dots,4).
\end{gathered}
\end{equation}
Thus $A_{11}=M=\pi_{K_0}A|_{K_0}$ is the compression, $A_{21}=P_{K_0^\perp}A|_{K_0}$ the coupling to
$K_0^\perp$, and $A_{12}=A_{22}=0$. For $d_1:=f_1+f_2$ and $d_2:=f_3+f_4$ one has
\begin{gather*}
Ad_1=d_1,\\
Ad_2=Md_2+A_{21}d_2=2f_3+3f_4+e_2=2d_2+(f_4+e_2),
\end{gather*}
because $A_{21}d_1=0$ and $A_{21}d_2=e_2$. Hence
\[
U:=D+AD=\operatorname{span}\{d_1,\ d_2,\ f_4+e_2\},\qquad \dim U=3 ,
\]
the three vectors being independent because $f_4+e_2\notin K_0$ while $d_1,d_2\in K_0$. Every
$\alpha d_1+\beta d_2+\gamma(f_4+e_2)\in K_0$ has $\gamma=0$, so
\[
K_0\cap U=K_0\cap(D+AD)=D .
\]
Put $V:=U^\perp$, so $q=\dim V=8-3=5$, and take $\mathcal L$ to be a matrix whose rows form a basis of
$V$; for $L_0=0$ this is a design with row space $V$.
Then $D=\ker H=K_0\cap V^\perp$, $\omega=\rank H=k-\dim D=2$, $c=k-\omega=2$ and $t=q-\omega=3$; the
kernel is exact because $K_0\cap V^\perp=K_0\cap U=D$. Moreover $AD\subseteq U=V^\perp$, so
$\ker H\subseteq\ker H_A$ and \eqref{eq:5a} holds. Finally
\[
V=U^\perp=\operatorname{span}\{e_1,\ e_3,\ e_4,\ e_6-e_5,\ e_8-e_7-e_2\},
\]
whence $V\cap K_0=\operatorname{span}\{e_6-e_5\}$ has dimension $1\ne\omega=2$: the design is not
$K_0$-split. And $Md_2=2f_3+3f_4=2d_2+f_4\notin D$, so $M(D)\not\subseteq D$.
\end{proof}

The split hypothesis is therefore a genuine restriction on the layer rather than a normalisation;
consequently every structural statement below is labelled \textbf{S}, and the only statement for the
unrestricted problem is Theorem~\ref{thm:ufibre}.

\subsection{The layer index is not monotone}\label{sec:layers:c}

At a fixed order $q$ the identity $c-t=k-q$ shows that $c$ and $t$ move together, so only one order
allows both layers to be compared. Take $k=4$: then $c=1$ forces $q=3+t$ and $c=2$ forces $q=2+t$, and
$q=3$ is the common value ($c=1,t=0$; $c=2,t=1$). For the instance $A,C$ of \cite{companion3}, \S VI,
with the compression replaced by $\operatorname{diag}(1,1,2,3)$:

\begin{thm}[non-monotonicity]\label{thm:nonmono}
At order $q=3$ the $c=2$ layer is feasible and the $c=1$ layer is not.
\end{thm}

\begin{proof}
\emph{$c=2$ side.} With $D=\operatorname{span}\{e_1+e_2,e_3\}$ and
$W=\operatorname{span}\{(0,0,1,1,0,0,0,0)\}$, the design $V=D^{\perp_{K_0}}\oplus W$ has $\dim V=3$,
$\omega=2$, satisfies \eqref{eq:5a} and exactness ($\dim(K_0\cap V^\perp)=2$), and the six nonzero maximal
$2\times2$ minors of the pencil are
\[
-(\lambda-3)(\lambda-1),\ (\lambda-3)(\lambda-1),\ \lambda-1,\ 1-\lambda,\ \lambda-3,\ 3-\lambda,
\]
whose gcd is $1$ (indeed $\tfrac12\bigl((\lambda-1)-(\lambda-3)\bigr)=1$). So \eqref{eq:5b} holds.

\emph{$c=1$ side.} At $q=3$ and $c=1$ one has $t=q-\omega=0$, so $\dim(V\cap K_0^\perp)=0$ and the design
is automatically split; Claim~\ref{clm:invariant} therefore applies. The invariant lines are the lines of
the eigenspace $E_1=\operatorname{span}\{e_1,e_2\}$ together with $\operatorname{span}\{e_3\}$ and
$\operatorname{span}\{e_4\}$. For the family $D=\operatorname{span}\{e_1+se_2\}$ one computes
\[
H=\begin{pmatrix}-s&1&0&0\\0&0&1&0\\0&0&0&1\end{pmatrix},\qquad \rank H=3,
\]
and the gcd of the nonzero maximal minors equals $(\lambda-1)(\lambda-2)(\lambda-3)$,
\textbf{independently of $s$}; the point at infinity $s=\infty$ gives the same gcd, and the two isolated
directions give $(\lambda-1)^2(\lambda-3)$ and $(\lambda-1)^2(\lambda-2)$. Hence no $c=1$ design of order
$3$ is feasible.
\end{proof}

\begin{remark}\label{rem:direction}
The direction of the failure is ``the larger layer is easier'', so no monotonicity argument can replace
the analysis of all layers.
\end{remark}

Consequently the analysis of one layer cannot replace the analysis of another at the same order, and
Sections~\ref{sec:fibres}--\ref{sec:decision} treat $c$ as given.

\section{The fixed-kernel fibre}\label{sec:fibres}

This section computes, for a fixed kernel direction $D$, the exact set of admissible row spaces in the two
problems. Part A treats the unrestricted problem \textbf{U}; Part B the split problem \textbf{S}. In both
cases the answer is a Grassmannian fibre with an explicit dimension and an explicit non-emptiness
condition.

Throughout this section $D\le K_0$ with $\dim D=c$ is fixed, $U_D:=D+AD$, and $V_0=\row L_0$. We write
$P_{K_0}$ and $P_{K_0^\perp}$ for the two orthogonal projections, and $M=P_{K_0}A|_{K_0}$,
$A_{21}=P_{K_0^\perp}A|_{K_0}$, so that $Ax=Mx+A_{21}x$ for $x\in K_0$ (the two summands lie in
complementary subspaces).

\subsection{The unrestricted problem \texorpdfstring{\textbf{U}}{U}}\label{sec:fibres:u}

\begin{thm}[unrestricted fibre]\label{thm:ufibre}
Let $D\le K_0$ with $\dim D=c$.
\begin{enumerate}[label=\arabic*.,leftmargin=*,itemsep=1pt,topsep=2pt]
\item A $q$-dimensional $V$ satisfies $V_0\subseteq V$, $V\subseteq U_D^\perp$ and $K_0\cap V^\perp=D$ if
and only if
\[
V_0\subseteq V\subseteq U_D^\perp,\qquad \dim V=q,\qquad K_0\cap(D+AD)=D .
\]
\item The condition $K_0\cap(D+AD)=D$ is \textbf{equivalent} to exactness of the kernel, i.e. to
$K_0\cap V^\perp=D$ for every $V$ with $V_0\subseteq V\subseteq U_D^\perp$.
\end{enumerate}
\end{thm}

\begin{proof}
(1) If $K_0\cap V^\perp=D$ and \eqref{eq:5a} holds then $U_D\subseteq V^\perp$ and hence
$K_0\cap(D+AD)\subseteq K_0\cap V^\perp=D$, while $D\subseteq D+AD$ gives the reverse inclusion.
Conversely if $K_0\cap(D+AD)=D$, then any $V$ with $V_0\subseteq V\subseteq U_D^\perp$ satisfies
$D\subseteq K_0\cap V^\perp\subseteq K_0\cap(D+AD)=D$. (2) is the same two lines.
\end{proof}

\begin{remark}\label{rem:ufibre}
For $c=2$ the fibre of \textbf{U} is thus
$\{V\in\Gr(q,U_D^\perp):V_0\subseteq V,\ K_0\cap V^\perp=D\}$, which is an open subset of a Grassmannian
once the exactness condition holds. We do not eliminate the fibre variable in \textbf{U}.
\end{remark}

The unrestricted fibre is thus an open subset of a Grassmannian once the exactness condition
$K_0\cap(D+AD)=D$ holds, and that condition is a property of $D$ alone.

\subsection{The split problem \texorpdfstring{\textbf{S}}{S}}\label{sec:fibres:s}

\begin{thm}[split fibre]\label{thm:splitfibre}
Let $D\in\mathcal I_c(M)$ with $\dim D=c$, and put
\[
T_D:=K_0^\perp\cap(A_{21}D)^\perp,\qquad W_0:=P_{K_0^\perp}(V_0),
\]
\[
w_0:=\dim W_0,\qquad r(D):=\rank(A_{21}|_D).
\]
Assume the two prescribed-row conditions
\begin{equation}\label{eq:splithyp}
\textbf{(H1)}\quad P_{K_0}(V_0)\subseteq D^{\perp_{K_0}},
\qquad
\textbf{(H2)}\quad W_0\subseteq T_D .
\end{equation}
Then the row spaces of the $K_0$-split designs with kernel $D$, prescribed rows $V_0$, and satisfying
\eqref{eq:5a} are exactly
\[
V=D^{\perp_{K_0}}\oplus W,\qquad W_0\subseteq W\subseteq T_D,\quad \dim W=t,
\]
and the fibre is
\begin{equation}\label{eq:splitfibre}
\begin{gathered}
\mathcal F_D^{\rm split}\cong\Gr\bigl(t-w_0,\ T_D/W_0\bigr),\\
\dim\mathcal F_D^{\rm split}=(t-w_0)\bigl((n-k)-r(D)-t\bigr),
\end{gathered}
\end{equation}
which is non-empty if and only if
\[
w_0\le t\le (n-k)-r(D).
\]
\end{thm}

\begin{proof}
(i) For a split design, $\dim(V\cap K_0)=\omega=k-c$ by Claim~\ref{clm:identities}(4). Every
$x\in V\cap K_0$ is orthogonal to $D\subseteq V^\perp$, and $\dim D^{\perp_{K_0}}=k-\dim D=k-c$; hence
$V\cap K_0=D^{\perp_{K_0}}$, and $V=(V\cap K_0)\oplus W$ with $W:=V\cap K_0^\perp$. Then
$K_0\cap V^\perp=D$: the inclusion $D\subseteq K_0\cap V^\perp$ holds because $D\perp D^{\perp_{K_0}}$ and
$D\perp W$, and the reverse inclusion follows from
$K_0\cap V^\perp\subseteq K_0\cap(D^{\perp_{K_0}})^\perp=D$.
(ii) Recall that \eqref{eq:5a} is $\ker H\subseteq\ker H_A$, i.e.\ $AD\perp V$: indeed
$\ker H_A=\{u\in K_0:Au\in V^\perp\}$. For $d\in D$, $\langle d,V\rangle=0$; writing $Ad=Md+A_{21}d$
with $Md\in D\subseteq V^\perp$ and $A_{21}d\in K_0^\perp$, the relation $\langle Ad,V\rangle=0$ is
equivalent to $W\perp A_{21}D$. Hence \eqref{eq:5a} holds if and only if $W\subseteq T_D$, i.e.\ if and
only if the fibre parameter lies in $T_D$.
(iii) $V_0\subseteq V$ splits into $P_{K_0}(V_0)\subseteq D^{\perp_{K_0}}$ and $W_0\subseteq W$, which are
(H1) and the lower inclusion of the fibre.
(iv) Since $A_{21}D\subseteq K_0^\perp$ and $\dim A_{21}D=r(D)$, its orthogonal complement inside
$K_0^\perp$ has dimension $(n-k)-r(D)$; thus $\dim T_D=(n-k)-r(D)$, and $W/W_0$ is a free
$(t-w_0)$-plane in $T_D/W_0$, a Grassmannian of dimension
$(t-w_0)(\dim T_D-t)=(t-w_0)((n-k)-r(D)-t)$.
\end{proof}

\begin{remark}[the prescribed rows are two conditions]\label{rem:presrows}
The condition on the base,
\begin{equation}\label{eq:bpres}
\begin{aligned}
\mathcal B_{\rm pres}&:=\bigl\{D\in\mathcal I_c(M):\ P_{K_0}(V_0)\subseteq D^{\perp_{K_0}}\bigr\}\\
&=\bigl\{D\in\mathcal I_c(M):\ D\perp P_{K_0}(V_0)\bigr\},
\end{aligned}
\end{equation}
is not implied by $W_0\subseteq W$; omitting it produces ``designs'' that cannot be completed to a legal
observer. It also removes entire components of the base: if $P_{K_0}(V_0)$ contains a vector with
non-zero inner product with every element of a given invariant $D$, that $D$ is excluded regardless of
the fibre.
\end{remark}

The split fibre is thus a Grassmannian of computable dimension over the field of definition of $D$, and
its two prescribed-row conditions are independent; this is the description that
Section~\ref{sec:reduction} compresses further.

\section{Image-space reduction}\label{sec:reduction}

This section proves the reduction on which the rest of the paper rests: for a split design with
$M(D)\subseteq D$, the pencil sees the fibre $W$ only through its image
$\Phi(W)\subseteq\varphi(T_D)\subseteq D^\circ$. Part A proves the reduction, Part B gives the exact
realisability conditions for a prescribed image, and Part C states the resulting criterion. Nothing in
the section uses $c=2$; the reader interested only in the general statement may read it together with
Proposition~\ref{prop:general}.

Fix $D\in\mathcal B_{\rm pres}$ and a split design $V=D^{\perp_{K_0}}\oplus W$ with
$W_0\subseteq W\subseteq T_D$.

\subsection{The reduction}\label{sec:reduction:a}

Write $H_K,(H_A)_K$ for the rows of $H,H_A$ carried by $D^{\perp_{K_0}}$, and put
$P_D(\lambda):=\lambda H_K-(H_A)_K$, $R_\lambda:=\row P_D(\lambda)$.

Let
\begin{equation}\label{eq:iotaN}
\iota_N:K_0^*\longrightarrow\mathbb R^{1\times k},\qquad \ell\longmapsto \ell N
\end{equation}
be the coordinate isomorphism induced by the basis $N$: the row vector $\ell N$ has entries
$\ell(Ne_j)$, $j=1,\dots,k$, and $\iota_N$ is a linear isomorphism because $N$ has rank $k$. We identify
$K_0^*$ with $\mathbb R^{1\times k}$ through $\iota_N$; from now on every row-space identity is understood
after this identification, and a subspace of $K_0^*$ and its image under $\iota_N$ are written with the
same symbol.

For $v\in E$ let $h_v,h^A_v\in K_0^*$ be the restrictions to $K_0$ of $u\mapsto\langle v,u\rangle$ and
$u\mapsto\langle v,Au\rangle$, and put $\Psi_\lambda(v):=\lambda h_v-h^A_v$. The $j$-th row of
$\lambda H-H_A$ is $\iota_N(\Psi_\lambda(v_j))$, where $v_j$ is the $j$-th row of $\mathcal L$; hence
$\row(\lambda H-H_A)=\iota_N(\Psi_\lambda(V))$.

For $w\in K_0^\perp$ define $\varphi_w\in K_0^*$ by $\varphi_w(u)=\langle w,A_{21}u\rangle$, and let
\[
\varphi:K_0^\perp\longrightarrow K_0^*,\qquad \varphi(w):=\varphi_w
\]
be the induced linear map. For $W\le K_0^\perp$ put
\begin{equation}\label{eq:Phi}
\Phi(W):=\varphi(W)=\{\varphi_w:\ w\in W\}\le D^\circ .
\end{equation}
The inclusion is an inclusion of subspaces, and $\Phi(W)$ is the image of the linear map $\varphi$, not
merely a set: $W\subseteq T_D$ means $W\perp A_{21}D$, hence $\varphi_w(d)=\langle w,A_{21}d\rangle=0$ for
all $d\in D$, so $\varphi(W)\subseteq D^\circ$.

\begin{thm}[image-space reduction]\label{thm:reduction}
For every $\lambda$,
\begin{equation}\label{eq:reduction}
\row\bigl(\lambda H-H_A\bigr)=R_\lambda+\Phi(W),
\end{equation}
the identity holding in $\mathbb R^{1\times k}$ under $\iota_N$; equivalently,
$\Psi_\lambda(V)=\Psi_\lambda(D^{\perp_{K_0}})+\Phi(W)$ in $K_0^*$. Consequently the pencil condition
depends on $W$ only through the subspace $\Phi(W)$, and
$\dim\Phi(W)\le\dim\varphi(T_D)\le\dim T_D=(n-k)-r(D)$. Moreover:
\begin{enumerate}[label=\arabic*.,start=3,leftmargin=*,itemsep=1pt,topsep=2pt]
\item $\row(\lambda H-H_A)\subseteq D^\circ$ for every $\lambda$, and $\dim D^\circ=k-c=\omega$;
\item hence $\rank(\lambda H-H_A)\le\omega$ automatically, and \eqref{eq:5b} is equivalent to
\[
\row\bigl(\lambda H-H_A\bigr)=D^\circ\quad\text{for all }\lambda,
\]
\[
\text{i.e.}\qquad R_\lambda+\Phi(W)=D^\circ\quad\forall\lambda.
\]
\end{enumerate}
\end{thm}

\begin{proof}
\begin{enumerate}[label=\arabic*.,leftmargin=*,itemsep=1pt,topsep=2pt]
\item For $w\in W\le K_0^\perp$ and $u\in K_0$: $h_w(u)=\langle w,u\rangle=0$, and
$h^A_w(u)=\langle w,Au\rangle=\langle w,Mu\rangle+\langle w,A_{21}u\rangle=0+\varphi_w(u)$; hence
$\Psi_\lambda(w)=-\varphi_w$, independently of $\lambda$.
\item For any linear $\Psi$ and any direct sum $V=P\oplus W$, $\Psi(V)=\Psi(P)+\Psi(W)$; here
$P=D^{\perp_{K_0}}$.
\end{enumerate}
Combining, $\Psi_\lambda(V)=\Psi_\lambda(P)+\Phi(W)$, and the sign is immaterial because a subspace
equals its negative; applying the isomorphism $\iota_N$ gives \eqref{eq:reduction}.
\begin{enumerate}[label=\arabic*.,start=3,leftmargin=*,itemsep=1pt,topsep=2pt]
\item For $v\in P=D^{\perp_{K_0}}$ and $d\in D$: $\langle v,d\rangle=0$, and
$\langle v,Ad\rangle=\langle v,Md\rangle+\langle v,A_{21}d\rangle=0+0$ because $Md\in D$ and
$A_{21}d\in K_0^\perp$; hence every $\Psi_\lambda(v)$ annihilates $D$, i.e.\ lies in $D^\circ$, and
$\Phi(W)\subseteq D^\circ$ by \eqref{eq:Phi}. Since $\dim D^\circ=k-c=\omega$, the rank bound follows.
\item A subspace of the $\omega$-dimensional $D^\circ$ has dimension $\omega$ if and only if it is
$D^\circ$.
\end{enumerate}
\end{proof}

\begin{remark}[what the reduction says]\label{rem:whatreduction}
The spectral condition does not see $W$, only $\Phi(W)$: if $\Phi(W_1)=\Phi(W_2)$ then the two designs
have identical pencils up to row operations. The design space is therefore quotiented by the fibres of
$\Phi$, and the relevant parameter lives in the image $\varphi(T_D)$, a space of dimension at most $n-k$.
For $c=2$ this is decisive (Section~\ref{sec:decision}).
\end{remark}

So the spectral condition is a condition on the image $\Phi(W)$ and not on $W$, and the fibre variable
has been replaced by a subspace of $\varphi(T_D)$, a space whose dimension is at most the codimension of
$K_0$.

\subsection{Exact realisability}\label{sec:reduction:b}

\begin{lem}[subspaces with a prescribed image]\label{lem:prescribed}
Let $f:X\to Y$ be linear, let $Z\le\operatorname{im}f$ and let $d\ge0$. There is a subspace $S\le X$
with $f(S)=Z$ and $\dim S=d$ if and only if
\begin{equation}\label{eq:prescribed}
\dim Z\le d\le\dim\ker f+\dim Z .
\end{equation}
\end{lem}

\begin{proof}
If $f(S)=Z$ then $\dim Z=\dim S-\dim(S\cap\ker f)\le d$, and $S\le f^{-1}(Z)$ gives
$d\le\dim f^{-1}(Z)=\dim\ker f+\dim Z$. Conversely assume \eqref{eq:prescribed}. Choose a complement
$S_0$ of $\ker f$ in $f^{-1}(Z)$; then $f|_{S_0}:S_0\to Z$ is an isomorphism, so $\dim S_0=\dim Z$.
Since $\dim f^{-1}(Z)=\dim\ker f+\dim Z\ge d$, extend $S_0$ inside $f^{-1}(Z)$ to a subspace $S$ of
dimension $d$. Then $f(S)=f(S_0)=Z$.
\end{proof}

\begin{thm}[realisability of a prescribed image]\label{thm:realize}
Assume (H2) of \eqref{eq:splithyp}, so that $T_D/W_0$ is defined. Let $\Phi_0:=\varphi(W_0)$ and let
$\bar\varphi:T_D/W_0\to D^\circ/\Phi_0$ be the map induced by
$\varphi|_{T_D}$. A subspace $\Phi\le D^\circ$ is realised by
some $W$ with $W_0\subseteq W\subseteq T_D$ and $\dim W=t$ if and only if
\begin{equation}\label{eq:realize}
\begin{aligned}
&\Phi_0\subseteq\Phi\subseteq\varphi(T_D),\\
&\dim\Phi-\dim\Phi_0\ \le\ t-w_0\\
&\qquad\le\ \dim\ker\bar\varphi+\dim\Phi-\dim\Phi_0 .
\end{aligned}
\end{equation}
\end{thm}

\begin{proof}
Subspaces $W$ with $W_0\subseteq W\subseteq T_D$ correspond bijectively to subspaces
$S=W/W_0\le T_D/W_0$, with $\dim S=\dim W-w_0$, and $\varphi(W)=\Phi$ if and only if
$\bar\varphi(S)=\Phi/\Phi_0$ (which presupposes $\Phi_0\subseteq\Phi$). Apply
Lemma~\ref{lem:prescribed} to $f=\bar\varphi$, $Z=\Phi/\Phi_0$ and $d=t-w_0$: an $S$ with
$\bar\varphi(S)=Z$ and $\dim S=d$ exists if and only if $Z\le\operatorname{im}\bar\varphi$, i.e.\
$\Phi\subseteq\varphi(T_D)$, and
$\dim Z\le d\le\dim\ker\bar\varphi+\dim Z$, which is \eqref{eq:realize}.
\end{proof}

Every subspace of $D^\circ$ between $\Phi_0$ and $\varphi(T_D)$ whose dimension lies in the stated interval
is therefore realised, and no other subspace is. Since this set is what the decision procedure of
Section~\ref{sec:decision} enumerates, we name it, writing $V_\varphi:=\varphi(T_D)$:
\begin{equation}\label{eq:ID}
\begin{gathered}
\mathcal I_D:=\bigl\{\Phi\le V_\varphi:\ \Phi_0\subseteq\Phi,\\
\dim\Phi-\dim\Phi_0\ \le\ t-w_0\\
\le\ \dim\ker\bar\varphi+\dim\Phi-\dim\Phi_0\bigr\}.
\end{gathered}
\end{equation}
with $\bar\varphi:T_D/W_0\to D^\circ/\Phi_0$ the map of Theorem~\ref{thm:realize}. By that theorem,
$\Phi\in\mathcal I_D$ if and only if $\Phi=\Phi(W)$ for some $W$ with $W_0\subseteq W\subseteq T_D$ and
$\dim W=t$; so $\mathcal I_D$ is exactly the set of images of the fibre, and it is over $\mathcal I_D$ that
the quantifier of Theorem~\ref{thm:reduced} runs. The set \eqref{eq:ID} is \emph{not} the whole union
$\bigcup_j\Gr(j,V_\varphi)$: the two dimension conditions confine $j=\dim\Phi$ to a finite range of values,
and $\Phi_0\subseteq\Phi$ confines the position of $\Phi$ inside each Grassmannian.
Section~\ref{sec:algorithm} describes how the algorithm enumerates $\mathcal I_D$ by finitely many charts.


\subsection{The reduced criterion}\label{sec:reduction:c}

\begin{thm}\label{thm:reduced}
Recall $P_D(\lambda)=\lambda H_K-(H_A)_K$ and $r_D(\lambda)=\dim R_\lambda$. Then
\begin{align}
&\exists W\ (W_0\subseteq W\subseteq T_D,\ \dim W=t)\ \text{satisfying \eqref{eq:5b}}\notag\\
&\iff\ \exists\Phi\in\mathcal I_D\ \text{(the realisable set \eqref{eq:ID}) such that}\notag\\
&\qquad\dim\bigl(\row P_D(\lambda)+\Phi\bigr)=\omega\quad\forall\lambda,
\label{eq:reduced}
\end{align}
and the latter is equivalent to
\begin{equation}\label{eq:reduced2}
\dim\Phi-\dim\bigl(\Phi\cap\row P_D(\lambda)\bigr)=\omega-r_D(\lambda)\quad\forall\lambda .
\end{equation}
\end{thm}

\begin{proof}
Immediate from Theorem~\ref{thm:reduction} and $\dim(X+Y)=\dim X+\dim Y-\dim(X\cap Y)$.
\end{proof}

\begin{remark}\label{rem:imagesize}
Theorems~\ref{thm:reduction}--\ref{thm:reduced} reduce the fibre problem to a quantifier over subspaces of
$D^\circ$. Two spaces matter and they must not be confused: $D^\circ$ has dimension $\omega$, which grows
with $k$, while the image $\varphi(T_D)$ has dimension at most
$\dim T_D=(n-k)-r(D)\le n-k\le 2p$, where $p$ is the number of output rows. So the effective parameter
does not range over the whole union $\bigcup_j\Gr(j,\varphi(T_D))$ but over the realisable set
$\mathcal I_D$ of \eqref{eq:ID}, which is contained in
\[
\bigcup_{j\le\min(t,\,2p)}\Gr\bigl(j,\ \varphi(T_D)\bigr),
\]
a family of dimension at most $\lfloor\min(t,2p)^2/4\rfloor\le p^2$, \textbf{independent of $c$}. In
particular the image parameter is a point, a projective line, or a plane only when
$\dim\varphi(T_D)\le2$; for larger output counts the parameter is a higher-dimensional Grassmannian. A
necessary condition for feasibility follows at once:
$\max_\lambda(\omega-r_D(\lambda))\le\dim\Phi\le n-k$, i.e. the rank deficiency of $P_D(\lambda)$ can
never exceed the codimension of $K_0$.
\end{remark}

\section{The base at \texorpdfstring{$c=2$}{c=2}}\label{sec:base}

This section describes the base of invariant two-planes at $c=2$ under simple spectrum over $\mathbb C$,
and then separates the part of the construction that is valid for every $c$ from the parts that are not.
Part A is the classification, Part B the structural layer.

Write $\mathcal I_2(M)$ for $\mathcal I_c(M)$ with $c=2$ (Section~\ref{sec:layers}).

\emph{Simple spectrum over $\mathbb C$.} Throughout the paper the hypothesis means that the eigenvalues
of $M$ \emph{in $\mathbb C$} are pairwise distinct, equivalently that $\chi_M$ is squarefree. It is not
the weaker and different condition that the \emph{real} eigenvalues are distinct: a real matrix with
$\chi_M(x)=(x^2+1)(x^2+2)$, for instance the companion matrix of that polynomial, has simple spectrum over
$\mathbb C$ although it has no real eigenvalue at all. What the enumeration below uses is squarefreeness,
i.e.\ the factorization of $\chi_M$ into distinct irreducible factors over $\mathbb Q$.

\subsection{Simple spectrum over \texorpdfstring{$\mathbb C$}{C}: the invariant-plane base}\label{sec:base:a}

\begin{thm}[base, simple spectrum over $\mathbb C$]\label{thm:base}
Assume (i) $c=2$ and (ii) $M$ has simple spectrum over $\mathbb C$, i.e.\ its eigenvalues in
$\mathbb C$ are pairwise distinct, equivalently $\chi_M$ is squarefree. Then $\mathcal I_2(M)$ is finite,
and every element is one of:
\begin{enumerate}[label=\arabic*.,leftmargin=*,itemsep=1pt,topsep=2pt]
\item \textbf{(branch R)} $\operatorname{span}\{v_\alpha,v_\beta\}$ for two real eigenvalues
$\alpha\ne\beta$ with eigenvectors $v_\alpha,v_\beta$;
\item \textbf{(branch C)} the real two-plane attached to a complex-conjugate pair. Let
$\mu=a+ib$, $b\ne0$, be an eigenvalue with eigenvector $z=u+iw$, $u,w$ real, and let $f\in\mathbb Q[x]$
be the minimal polynomial of $\mu$ over $\mathbb Q$, of degree $d=\deg f\le k$. Then
\[
Mu=au-bw,\qquad Mw=bu+aw,
\]
so $M(D)\subseteq D$ for $D:=\operatorname{span}_{\mathbb R}\{u,w\}$. The plane $D$ is defined over the
real field $\mathbb Q(\mu+\bar\mu,\mu\bar\mu)\subseteq\mathbb Q(\mu,\bar\mu)$, of degree at most
$d^2\le k^2$ over $\mathbb Q$; it is \emph{not} in general defined over $\mathbb Q$. The two statements
$D=\ker f(M)$ and ``$D$ is defined over $\mathbb Q$'' both hold exactly when $d=2$, i.e.\ when the
conjugate orbit of $\mu$ is $\{\mu,\bar\mu\}$. For $d>2$ the kernel $\ker f(M)$ is the span of the real
forms of the eigenspaces of all the conjugates of $\mu$, and has dimension $d$, not $2$.%
\footnote{The distinction is not vacuous. For $M$ the companion matrix of $x^4+1$ and
$\mu=e^{i\pi/4}$ one has $d=4$, so $\chi_M=f$ is irreducible and
$\ker f(M)=\ker(M^4+I)=\mathbb R^4$, whereas $D$ is two-dimensional; the Pl\"ucker coordinates of $D$
are proportional to $(1,\sqrt2,1,1,\sqrt2,1)$, so $D$ is defined over $\mathbb Q(\sqrt2)$ and not over
$\mathbb Q$.}
\end{enumerate}
The total number of planes is $\binom{k_{\mathbb R}}{2}+\#\{\text{conjugate pairs}\}\le k^2$, where
$k_{\mathbb R}$ is the number of real eigenvalues.
\end{thm}

\begin{proof}
For simple spectrum over $\mathbb C$, $M$ is diagonalisable over $\mathbb C$ and each complex eigenspace
is a line. An invariant real two-plane either contains two real eigenvectors (branch R), or is spanned by
the real and imaginary parts of a complex eigenvector (branch C), or is $\mathbb R^2$ for a
two-dimensional $M$ with a non-real pair, which is the case $k=2$ of branch C. Nothing else is invariant,
because the restriction of $M$ to $D$ has a characteristic polynomial of degree two dividing the
squarefree $\chi_M$.

For branch C the two identities are explicit. If $Mz=\mu z$ with $z=u+iw$ and $\mu=a+ib$, comparing real
and imaginary parts gives $Mu=au-bw$ and $Mw=bu+aw$, hence $M(D)\subseteq D$ for
$D=\operatorname{span}_{\mathbb R}\{u,w\}$. The vectors $u,w$ are independent over $\mathbb R$, since a
real eigenvector of $M$ would have the real eigenvalue $a$.

Let $f\in\mathbb Q[x]$ be the minimal polynomial of $\mu$ and $d=\deg f\le k$. Since $\chi_M$ is
squarefree, $f$ occurs in $\chi_M$ exactly once, so the eigenspaces $E_\nu$ of the roots $\nu$ of $f$
are one-dimensional and independent over $\mathbb C$, and
\[
\ker_{\mathbb C}f(M)=\bigoplus_{\nu:\,f(\nu)=0}E_\nu .
\]
As $f(M)$ is real, $\ker_{\mathbb R}f(M)=\ker_{\mathbb C}f(M)\cap\mathbb R^k$ has real dimension $d$; it
is the direct sum of the planes
$D_\nu=\operatorname{span}_{\mathbb R}\{\operatorname{Re}z_\nu,\operatorname{Im}z_\nu\}$ over the
conjugate pairs among the roots of $f$ and of the real eigenlines $v_\nu$ over its real roots. Hence
$D=D_\mu\subseteq\ker f(M)$, with equality exactly when $f$ has no root besides $\mu$ and $\bar\mu$,
i.e.\ when $d=2$.

For the field, choose $z$ with coordinates in $\mathbb Q(\mu)$, solving $(\mu I-M)z=0$ over
$\mathbb Q(\mu)$, where the rank is $k-1$ because $\mu$ is a simple root of $\chi_M$; scaling $z$ by a
non-zero complex number, in particular by an element of $\mathbb Q(\mu)$, does not change $D$. The polynomial
$g(x):=x^2-(\mu+\bar\mu)x+\mu\bar\mu=(x-\mu)(x-\bar\mu)$ has coefficients in the real field
$L:=\mathbb Q(\mu+\bar\mu,\mu\bar\mu)\subseteq\mathbb Q(\mu,\bar\mu)$ and the two simple roots
$\mu,\bar\mu$, so $\ker_{\mathbb R}g(M)=D$: the complex kernel is $E_\mu\oplus E_{\bar\mu}$ and $D$ is its
space of real points. As $g(M)$ is an $L$-matrix, Cramer's rule over $L$ gives a basis of $D$ in
$L^k$, so $D$ is defined over $L$. Conversely, if $D$ is defined over a field $K$, then $M|_D$ is a
$K$-rational $2\times2$ matrix in a $K$-basis of $D$, so its trace $\mu+\bar\mu$ and determinant
$\mu\bar\mu$ lie in $K$; thus
$L$ is the smallest field of definition of $D$. Finally $L\subseteq\mathbb Q(\mu,\bar\mu)$, because
$\mu+\bar\mu$ and $\mu\bar\mu$ are polynomial expressions in $\mu$ and $\bar\mu$; hence
$[L:\mathbb Q]\le[\mathbb Q(\mu,\bar\mu):\mathbb Q]\le d^2\le k^2$, the middle bound being the degree of a
compositum of two extensions of degree $d$. Moreover
$L=\mathbb Q$ exactly when $d=2$: if $d=2$ then $g\in\mathbb Q[x]$; conversely $L=\mathbb Q$ makes $\mu$
a root of the quadratic $g\in\mathbb Q[x]$, whose roots are non-real, so $d=2$.

The count is immediate.
\end{proof}

\begin{remark}[what is \emph{not} claimed]\label{rem:nonsimple}
With repeated eigenvalues $\mathcal I_2(M)$ \emph{may} be positive-dimensional, and when it is, the
dimension and the geometry of the family are governed by the Jordan structure of $M$, not by the multiset
of eigenvalues alone. For instance, an eigenspace $K_\lambda$ of dimension at least three contributes the
family $\Gr(2,K_\lambda)$; two eigenspaces $K_\lambda,K_\mu$ of dimension at least two contribute
$\mathbb P(K_\lambda)\times\mathbb P(K_\mu)$; and a repeated irreducible quadratic factor of multiplicity
$g$ contributes at least the $M$-stable complex lines in $\ker f(M)$, a $\mathbb{CP}^{g-1}$-family of
real dimension $2(g-1)$. Positive dimension is not automatic, however: a repeated eigenvalue carried by a
single Jordan block contributes only the members of its invariant flag, hence finitely many planes. We
make \textbf{no} claim of a rational parametrisation of the general invariant-plane variety; the
simple-spectrum-over-$\mathbb C$ statement above is what the algorithm uses.
\end{remark}

\begin{table*}[!t]
\centering
\caption{The invariant-plane classification of Theorem~\ref{thm:base} at $c=2$ under simple spectrum over
$\mathbb C$: the two branches, the field over which each plane is defined (Claim~\ref{clm:arith}), and
the number of planes. Here $k_{\mathbb R}$ is the number of real eigenvalues of $M$ and $k=\dim K_0$.}
\label{tab:planes}
\small
\begin{tabular}{@{}l>{\raggedright\arraybackslash}p{0.44\textwidth}>{\raggedright\arraybackslash}p{0.24\textwidth}l@{}}
\toprule
Branch & Plane $D$ & Defined over & Number \\
\midrule
R & $\operatorname{span}\{v_\alpha,v_\beta\}$, two real eigenvalues $\alpha\ne\beta$
& $\mathbb Q(\alpha,\beta)$, degree $\le k^2$ & $\binom{k_{\mathbb R}}{2}$ \\
C & $\operatorname{span}\{u,w\}$ from $z=u+iw$ of $\mu=a+ib$, $b\ne0$; equal to
$\ker f(M)$ only when $\deg f=2$ & $\mathbb Q(\mu+\bar\mu,\mu\bar\mu)\subseteq\mathbb Q(\mu,\bar\mu)$,
degree $\le k^2$ ($=\mathbb Q$ exactly when $\deg f=2$) & $\#\{\text{conjugate pairs}\}$ \\
\midrule
Total & & & $\le k^2$ \\
\bottomrule
\end{tabular}
\end{table*}

Under simple spectrum over $\mathbb C$ the base is therefore finite and explicit, which is what makes the
enumeration of Section~\ref{sec:algorithm} terminate, and --- with the output dimension fixed ---
polynomial in the bit size of the input (Theorem~\ref{thm:complexity}); Remark~\ref{rem:nonsimple}
records what is not claimed without it.

\subsection{The general \texorpdfstring{$c$}{c}}\label{sec:base:b}

\begin{prop}[the structural layer at general $c$]\label{prop:general}
The following statements hold for every $c\ge0$ (the case $c=2$ is the one used in
Section~\ref{sec:decision}):
\begin{enumerate}[label=(\roman*),leftmargin=*,itemsep=1pt,topsep=2pt]
\item the fixed-kernel fibre descriptions of Theorems~\ref{thm:ufibre} and~\ref{thm:splitfibre}: the base
is $\mathcal B_{\rm pres}\subseteq\mathcal I_c(M)$, the fibre is $\Gr(t-w_0,T_D/W_0)$ with
$T_D=K_0^\perp\cap(A_{21}D)^\perp$;
\item the image-space reduction of Theorem~\ref{thm:reduction}:
$\row(\lambda H-H_A)=R_\lambda+\Phi(W)$ with $\Phi(W)\subseteq D^\circ$ and $\dim D^\circ=k-c=\omega$;
\item the realisability criterion of Lemma~\ref{lem:prescribed} and Theorem~\ref{thm:realize}, and the
reduced criterion of Theorem~\ref{thm:reduced}, unchanged; the effective image parameter ranges over the
realisable set $\mathcal I_D$ of \eqref{eq:ID}, which is contained in
$\bigcup_{j\le\min(t,2p)}\Gr(j,\varphi(T_D))$ and has dimension at most
$\lfloor\min(t,2p)^2/4\rfloor\le p^2$.
\end{enumerate}
Two further statements are \emph{not} covered by the proposition. First, the finite enumeration of the
base is a $c=2$ statement: Theorem~\ref{thm:base} enumerates $\mathcal I_2(M)$ under simple spectrum over
$\mathbb C$, whereas for $c\ge3$ the base $\mathcal I_c(M)$ is known only through its defining equations
(Remark~\ref{rem:open}). Second, the correctness and the polynomial-time bound are simple-spectrum-only:
Theorem~\ref{thm:correct} and Theorem~\ref{thm:complexity} are statements about the algorithm, they
consume the enumeration of Theorem~\ref{thm:base}, and they are \emph{not} generalised here.
\end{prop}

\begin{proof}
The proofs of Theorems~\ref{thm:ufibre}--\ref{thm:reduced} use only $\dim D=c$ and, where a base is
involved, $M(D)\subseteq D$; the numerical value of the layer index enters nowhere in them. What $c=2$
buys is the finiteness of $\mathcal I_2(M)$ under simple spectrum over $\mathbb C$
(Theorem~\ref{thm:base}), and that finiteness is an input to the algorithm, not a consequence of
clauses (i)--(iii). Theorems~\ref{thm:correct} and~\ref{thm:complexity} therefore inherit both the $c=2$
restriction and the simple-spectrum hypothesis, and the proposition asserts nothing about them.
\end{proof}

\begin{remark}\label{rem:annihilator}
Note that $\dim D^\circ=\omega$, not $c$: since $c=k-\omega$, the annihilator of the $c$-dimensional
subspace $D$ in the $k$-dimensional space $K_0^*$ has dimension $k-c=\omega$. The annihilator of
$V\cap K_0$ is a different space: in the split case it has dimension $k-\dim(V\cap K_0)=c$, and it is
$D^\circ$, not $(V\cap K_0)^\circ$, that contains the pencil. For $c=2$ and $k=4$ the two dimensions
coincide, which is why the small cases are misleading; in general $D^\circ$ is large and only the image
$\varphi(T_D)$ is small.
\end{remark}

\begin{remark}[what remains open]\label{rem:open}
For general $c$ three obstructions persist: the enumeration of $\mathcal I_c(M)$ (exponential when $c$
varies), the choice of $\Phi$ (a Grassmannian of dimension up to $p^2$), and the absence of a rational
parametrisation of the base in the non-simple case. Two questions seem most promising: (a) whether the
quantifier over $\Phi$ can be replaced by a rank or separation condition on the induced maps of
$T_D/W_0$; (b) whether the layer index can be traded for $\dim U$ with $U=\pi_1(W)$, since
$c=k-\dim U$. We do not know the answer to either.
\end{remark}

The reduction of Section~\ref{sec:reduction} therefore does not use $c=2$; what $c=2$ buys is a finite
base (Theorem~\ref{thm:base}, under simple spectrum over $\mathbb C$), while the smallness of the image
parameter in Section~\ref{sec:decision} comes from the bound on $n-k$, not from the smallness of $c$
(Remark~\ref{rem:boundaries}).

\section{Decision and its cost}\label{sec:decision}

This section assembles the decision procedure at $c=2$: Part A shows that admissibility is automatic on
an chart, Part B states the algorithm, Part C proves it correct, and Part D bounds its bit
cost.

\subsection{Admissibility is automatic on every chart}\label{sec:admissible}

A \emph{chart} is a rational parametrization $s\mapsto\Phi(s)$ of a family of image spaces inside
$\bigcup_j\Gr(j,V_\varphi)$, together with the pencil $(\lambda H-H_A)(s)$ that it induces. The parameter
$s$ ranges over a nonempty Zariski-open set $S=\{s\in\mathbb R^N:\ \delta(s)\neq0\}$ of $\mathbb R^N$; the
entries of $H$ and $H_A$ are rational functions of $s$ whose denominators divide $\delta$, and $\delta$ is
the \emph{chart denominator}. Clearing denominators, i.e.\ passing to $\delta^{\,e}(\lambda H-H_A)$ with $e$
large enough, multiplies every maximal minor by the same element $\delta^{\,e\omega}$ of the function field
and therefore changes no gcd in $\lambda$. A chart is a \emph{standard affine chart} of $\Gr(j,V_\varphi)$
when $s$ is the matrix $X$ of a linear map $U\to U'$ with $V_\varphi=U\oplus U'$ and $\dim U=j$, so that
$\Phi$ is the graph of $X$; on such a chart the dependence on $s$ is affine
(Claim~\ref{clm:degrees}), so $\delta=1$ there, and a denominator is needed only when the parametrization is
normalized by the eigenvector data of $D$ in branch~R of Theorem~\ref{thm:base}. A point $s\in S$ is
\emph{admissible} when some maximal minor of the cleared pencil has $\lambda$-degree $\omega$ at $s$; since
$\deg_\lambda\le\omega$ always, this asks that some maximal minor keep the largest $\lambda$-degree it can
have. The coefficient field is $K=L_D$, the field over which $D$ is defined in
Theorem~\ref{thm:base}: $L_D=\mathbb Q(\alpha,\beta)$ in branch~R and
$L_D=\mathbb Q(\mu+\bar\mu,\mu\bar\mu)$ in branch~C, of degree $d_D\le k^2$ in either branch; every
parameter space below is an affine space, hence irreducible, so a chart has a
single \emph{function field} $K(s)$, the field of rational functions in its parameters.

Recall from \cite{companion3} that the generic-gcd criterion rests on two facts in this situation: at every
point of a chart the degree of the generic $\lambda$-gcd bounds the degree of the specialised gcd from
below, and a constant generic gcd forces the chart to contain a real feasible point, because the real points
are Zariski dense. Lemma~\ref{lem:genericgcd} states and proves both facts for the charts at hand;
Claim~\ref{clm:admissible} first shows that every point of every chart is admissible, so no point ever has
to be excluded.

\begin{claim}[fixed leading coefficient]\label{clm:admissible}
On every chart of a split design, $H=\mathcal L N$ does not depend on the chart parameter: its $W$-rows
vanish and its $D$-rows are the $\omega$ independent rows of $H_K$, so $\rank H=\omega$ identically.
Let $I\subseteq\{1,\dots,q\}$ and $J\subseteq\{1,\dots,k\}$ be index sets with $|I|=|J|=\omega$ and
$\det H_{I,J}\neq0$, and let
\[
m_{I,J}(\lambda,s):=\det\bigl((\lambda H-H_A(s))_{I,J}\bigr)
\]
be the corresponding maximal minor of the pencil. Then
\begin{equation}\label{eq:leadcoef}
[\lambda^{\omega}]\,m_{I,J}(\lambda,s)=\det H_{I,J}\neq0
\qquad\text{for every }s,
\end{equation}
independently of the chart parameter $s$ and of the choice of the index sets. Hence
$\deg_\lambda m_{I,J}(\cdot,s)=\omega$ at every point of
every chart: no specialisation of the parameters can lower the $\lambda$-degree of this minor, some maximal
minor keeps $\lambda$-degree $\omega$ everywhere, and every real point of every chart is admissible.
\end{claim}

\begin{proof}
The $W$-rows of $H$ vanish (Theorem~\ref{thm:reduction}, step~1) and the $D$-rows have rank
$\dim D^{\perp_{K_0}}=k-c=\omega$; the entries of $H$ are therefore constants, independent of $W$ and of the
chart parameter, and $\rank H=\omega$ identically. Index sets with $\det H_{I,J}\neq0$ exist because
$\rank H=\omega$. The submatrix of the pencil is $(\lambda H-H_A(s))_{I,J}=\lambda H_{I,J}-(H_A(s))_{I,J}$; by the Leibniz
formula its determinant is the signed sum, over the permutations $\sigma$ of $\{1,\dots,\omega\}$, of the
products $\prod_{a=1}^{\omega}\bigl(\lambda H_{I_a,J_{\sigma(a)}}-(H_A(s))_{I_a,J_{\sigma(a)}}\bigr)$. The
coefficient of $\lambda^{\omega}$ arises only from choosing the first summand in each of the $\omega$
factors; it is
$\sum_{\sigma}\operatorname{sgn}(\sigma)\prod_{a=1}^{\omega}H_{I_a,J_{\sigma(a)}}=\det H_{I,J}$, a non-zero
constant and in particular independent of $s$, while no entry of $H_A$ can contribute at degree $\omega$.
Therefore
$\deg_\lambda m_{I,J}(\cdot,s)=\omega$ for every $s$, and a specialisation $s_0$ that lowered this degree
would have to annihilate $\det H_{I,J}$, which is impossible.
\end{proof}

\begin{lem}[specialization and generic $\lambda$-gcd]\label{lem:genericgcd}
Let $M_1,\dots,M_r\in K[s,\lambda]$, $r\ge1$, be the maximal minors of a cleared pencil of a chart,
polynomial in $s=(s_1,\dots,s_N)$ and in $\lambda$, and assume:
\begin{enumerate}[label=\textbf{(H\arabic*)},leftmargin=*,itemsep=1pt,topsep=2pt]
\item $K$ is a number field embedded in $\mathbb R$, so $K\subseteq\mathbb R$ is infinite and of
characteristic $0$, and the coordinate ring $K[s]=K[s_1,\dots,s_N]$ is a unique factorisation domain with
fraction field $K(s)$;
\item the chart's real parameter space is a non-empty real affine space, written $\mathbb R^N$ after an
affine change of coordinates, and the chart denominators do not vanish on it: the chart proper is the
non-empty Zariski-open set $S=\{s\in\mathbb R^N:\ \delta(s)\neq0\}$ for some $0\neq\delta\in K[s]$;
\item some $M_{i_0}$ has $\lambda$-degree exactly $d_0\ge0$, and its $\lambda$-leading coefficient is a unit
of $K[s,1/\delta]$, hence non-vanishing on $S$; a non-zero constant of $K$ suffices, and that is the case
of Claim~\ref{clm:admissible};
\item the density used is that of the $K$-rational, equivalently real, points of a non-empty real affine
space: a non-zero polynomial in $N$ variables over $K$ vanishes neither on all of $\mathbb R^N$ nor on all
of $K^N$.
\end{enumerate}
Let $g\in K(s)[\lambda]$ be the monic gcd of $M_1,\dots,M_r$ over the function field $K(s)$, the field of
rational functions in the chart parameters over the coefficient field $K$; call $g$ the \emph{generic
$\lambda$-gcd} of the chart. Then:
\begin{enumerate}[label=(\arabic*),leftmargin=*,itemsep=1pt,topsep=2pt]
\item $g$ lies in $K[s,1/\delta][\lambda]$ and is monic in $\lambda$ (in $K[s][\lambda]$ when the leading
coefficient in \textbf{(H3)} is a constant); its $\lambda$-degree is therefore the same at every point of
$S$, namely $\deg_\lambda g$, and $g(s,\cdot)$ divides every $M_i(s,\cdot)$ in $\mathbb R[\lambda]$ (in
$K[\lambda]$ when $s$ is $K$-rational);
\item if $\deg_\lambda g\ge1$, then
$\deg_\lambda\gcd\bigl(M_1(s,\cdot),\dots,M_r(s,\cdot)\bigr)\ge\deg_\lambda g\ge1$ for every $s\in S$, so
no point of the chart, admissible or not, has a constant $\lambda$-gcd;
\item if $\deg_\lambda g=0$, the set of $s\in S$ whose specialised gcd is not constant is contained in a
proper Zariski-closed subset of $S$; since the real points of a nonempty Zariski-open subset of
$\mathbb R^N$ are Zariski dense, $S$ contains a real point at which the specialised gcd is $1$.
\end{enumerate}
Consequently the generic $\lambda$-gcd has positive degree if and only if every point of $S$ has positive
specialised degree, and it has degree $0$ if and only if some real point of $S$ has constant specialised
gcd; when every point of $S$ is admissible, as on the charts of Claim~\ref{clm:admissible}, the same
statements hold with ``admissible specialisation'' in place of ``point of $S$''.
\end{lem}

\begin{proof}
Throughout, $A$ denotes one of the two rings $K[s]$ and $K[s,1/\delta]$; both are unique factorisation
domains with fraction field $K(s)$, and both contain the coefficients of all $M_i$. The proof needs no
localisation at points of $S$ and no regularity of $g$ at real points: the membership
$g\in K[s,1/\delta][\lambda]$ is deduced from a statement about monic divisors.

\emph{Step 0 (monic divisors over a UFD).} Let $B$ be a unique factorisation domain with fraction field
$F$, let $f\in B[\lambda]$ be monic, and let $h\in F[\lambda]$ be monic with $h\mid f$ in $F[\lambda]$. Then
$h\in B[\lambda]$. Indeed, over an algebraic closure $\bar F$ write
$f=\prod_{a=1}^{m}(\lambda-\alpha_a)$ and $h=\prod_{b=1}^{e}(\lambda-\beta_b)$; since $h\mid f$, every
$\beta_b$ equals some $\alpha_a$. A root of the monic $f\in B[\lambda]$ satisfies a monic equation with
coefficients in $B$, hence is integral over $B$, and the elements of $\bar F$ integral over $B$ form a
subring of $\bar F$ \cite[Prop.~5.1]{am69}; the coefficients of $h$, being the elementary symmetric
functions of the $\beta_b$, are therefore integral over $B$. They lie in $F$, and a unique factorisation
domain is integrally closed in its fraction field: if $x=a/b$ with $a,b\in B$ coprime satisfies
$x^n+c_{n-1}x^{n-1}+\dots+c_0=0$, then $b\mid a^n$, so $b$ is a unit. Hence $h\in B[\lambda]$: this is the
standard ``monic divisor of a monic polynomial over a UFD'' statement, here obtained by the root-coefficient
recursion (equivalently, by Gauss' lemma).

\emph{Step 1 (a monic reference minor).} By \textbf{(H3)} the minor $M_{i_0}$ has $\lambda$-degree $d_0$ and
$\lambda$-leading coefficient $u\in K[s,1/\delta]^{\times}$. Put
$\tilde M_0:=u^{-1}M_{i_0}\in K[s,1/\delta][\lambda]$; it is monic in $\lambda$ of degree $d_0$. If the
leading coefficient in \textbf{(H3)} is a non-zero constant $c\in K^{\times}$, then
$\tilde M_0=c^{-1}M_{i_0}\in K[s][\lambda]$. In the application $M_{i_0}=\delta^{e}m_{I,J}$ and
$u=\delta^{e}\det H_{I,J}$ with $\det H_{I,J}\in K^{\times}$, so $\tilde M_0=m_{I,J}/\det H_{I,J}$
(Claim~\ref{clm:admissible}, \eqref{eq:leadcoef}).

\emph{Step 2 (the generic gcd has coefficients in the chart ring).} By definition of the gcd, $g$ divides
$M_{i_0}$ in $K(s)[\lambda]$, hence divides the monic $\tilde M_0$ there. Step~0 with
$B=K[s,1/\delta]$, $f=\tilde M_0$ and $h=g$ gives $g\in K[s,1/\delta][\lambda]$; in the constant case of
Step~1 the same argument with $B=K[s]$ gives $g\in K[s][\lambda]$. (If $d_0=0$ then $\tilde M_0=1$ and
$g=1$.)

\emph{Step 3 (division).} For each $i$ we have $g\mid M_i$ in $K(s)[\lambda]$ by definition of the gcd, and
$g$ is monic with coefficients in $A$ by Step~2. Division by a monic polynomial never divides by a non-unit,
so the division algorithm runs inside $A[\lambda]$ and yields $M_i=g\,q_i$ with $q_i\in A[\lambda]$.

\emph{Step 4 (specialisation).} Let $s\in S$. Since $\delta(s)\neq0$, evaluation at $s$ is a ring
homomorphism $A\to\mathbb R$, with values in $K$ when $s$ is $K$-rational, so
$M_i(s,\lambda)=g(s,\lambda)q_i(s,\lambda)$ in $\mathbb R[\lambda]$. The specialised
polynomial $g(s,\lambda)$ is monic of degree $\deg_\lambda g$, because the leading coefficient $1$
specialises to $1$; it is therefore a common divisor of the specialised minors of that degree, and
$\deg_\lambda\gcd\bigl(M_1(s,\cdot),\dots,M_r(s,\cdot)\bigr)\ge\deg_\lambda g$. The bound thus holds at
\emph{every} point of $S$: the generic degree is a lower bound for every specialisation.

\emph{Step 5 (the degree-zero branch).} If $\deg_\lambda g=0$ then $g=1$, and the B\'ezout identity in
$K(s)[\lambda]$ gives $a_i\in K(s)[\lambda]$ with $\sum_i a_iM_i=1$. Clearing denominators gives
$b_i\in K[s,\lambda]$ and $0\neq\Delta\in K[s]$ with $\sum_i b_iM_i=\Delta$. At a point $s\in S$ with
$\Delta(s)\neq0$ the specialised minors have no common zero, hence gcd $1$; the exceptional set
$S\cap\{\Delta=0\}$ is therefore contained in the proper Zariski-closed set $\{\Delta=0\}$. By \textbf{(H4)}
the polynomial $\delta\Delta\in K[s]$ is not identically zero and vanishes neither on all of
$\mathbb R^N$ nor on all of $K^N$, so there is a point $s\in S$, rational over $K$ if desired, at which
$\Delta(s)\neq0$ and the specialised gcd is $1$. This is also where \textbf{(H4)} enters: a non-zero
polynomial vanishes on no non-empty Euclidean-open subset of $\mathbb R^N$, and $S$ contains one.

Statement~(1) is Steps 2--4, statement~(2) is Step~4, and statement~(3) is Step~5. The closing sentence
follows: if some point of $S$ has constant specialised gcd, Step~4 gives $\deg_\lambda g=0$, and if
$\deg_\lambda g=0$ then Step~5 exhibits a point of $S$ with specialised gcd $1$.
\end{proof}

Admissibility being automatic, and the generic-gcd test being exact on every chart of $\mathcal I_D$ by
Lemma~\ref{lem:genericgcd}, the algorithm can be stated in Part B without a certificate per point: a chart
contributes a feasible point exactly when its generic gcd is constant.

\subsection{The algorithm}\label{sec:algorithm}

\begin{defn}[image charts]\label{defn:charts}
Let $j=\dim\Phi$. Every $\Phi\in\mathcal I_D$ has
\begin{equation}\label{eq:jrange}
\begin{gathered}
\max\bigl\{\dim\Phi_0,\ t-w_0+\dim\Phi_0-\dim\ker\bar\varphi\bigr\}\ \le\ j\\
\le\ \min\bigl\{t-w_0+\dim\Phi_0,\ \dim V_\varphi\bigr\},
\end{gathered}
\end{equation}
and conversely, for every integer $j$ in this range, the set
$\{\Phi\in\Gr(j,V_\varphi):\ \Phi_0\subseteq\Phi\}$ is nonempty and is precisely the set of elements of
$\mathcal I_D$ of dimension $j$. The algorithm enumerates $\mathcal I_D$ through a finite chart cover: for
each $j$ in \eqref{eq:jrange} and each standard affine chart of $\Gr(j,V_\varphi)$ --- with parameters the
entries of a linear map $X:U\to U'$, $\dim U=j$, $V_\varphi=U\oplus U'$, so that $\Phi$ is the graph of $X$
--- the condition $\Phi_0\subseteq\Phi$ is a system of linear equations in $X$, so the realisable points of
that chart form an affine subspace $S\subseteq\mathbb R^{j(\dim V_\varphi-j)}$, possibly empty. These $S$
are the charts of Section~\ref{sec:admissible}; each carries at most $j(\dim V_\varphi-j)\le p^2$
parameters. If $t<w_0$ then $\mathcal I_D=\emptyset$ and $D$ is discarded, in agreement with the
non-emptiness condition of Theorem~\ref{thm:splitfibre}. The two endpoints of \eqref{eq:jrange} are part of
the same rule and need no separate treatment: $j=0$ is the single chart $\Phi=\{0\}$, realisable exactly
when $\Phi_0=0$ and the lower inequality of \eqref{eq:ID} holds, and $j=\dim V_\varphi$ is the single chart
$\Phi=V_\varphi$, which coincides with $D^\circ$ only in the boundary case $\dim T_D=\omega$. When
$\dim\Phi_0\ge1$ no one-dimensional family occurs; the line case occurs only for $\Phi_0=0$ and
$\dim V_\varphi\ge2$.
\end{defn}

The enumeration is finite and explicit, so what remains is to say on which charts the answer is decided
and at what cost.

\begin{algorithm}[H]
\caption{Decision at $c=2$, split, simple spectrum over $\mathbb C$.}
\label{alg:decide}
\begin{algorithmic}[1]
\State compute $K_0$, $N$, $M$, $A_{21}$, $\omega$, $c$, $t$;
\State enumerate $\mathcal I_2(M)$ by Theorem~\ref{thm:base}, keeping only
$D\in\mathcal B_{\rm pres}$;
\State for each $D$ compute $r(D)$, $T_D$, $W_0$, $\Phi_0$, and check $W_0\subseteq T_D$ and
$t\ge w_0$; discard $D$ otherwise;
\State enumerate the charts of $\mathcal I_D$: the finite range \eqref{eq:jrange} of $j=\dim\Phi$, and for
each $j$ the standard affine charts of $\Gr(j,V_\varphi)$ cut out by the linear condition
$\Phi_0\subseteq\Phi$ (Definition~\ref{defn:charts});
\State for each nonempty chart of $\mathcal I_D$, clear denominators and compute the gcd of the
$\omega\times\omega$ maximal minors over $K(s)[\lambda]$, $K=L_D$ the field of definition of $D$; if it is a non-zero
constant, return \emph{feasible};
\State if every chart of $\mathcal I_D$ gives a non-constant gcd, return \emph{infeasible}.

\end{algorithmic}
\end{algorithm}

\subsection{Correctness}\label{sec:correct}

\begin{claim}[exact yes/no decision]\label{clm:twovalued}
Assume (i) $c=2$, (ii) simple spectrum over $\mathbb C$, and (iii) the split condition, and let $S$ be a
nonempty chart of
$\mathcal I_D$ (Definition~\ref{defn:charts}). Then every point of $S$ is admissible
(Claim~\ref{clm:admissible}), so Lemma~\ref{lem:genericgcd} applies, with $M_1,\dots,M_r$ the maximal
minors of the cleared pencil in the chart parameter. Hence the generic $\lambda$-gcd of $S$ is a non-zero
constant if and only if $S$ contains a real point $s$ whose image $\Phi(s)\in\mathcal I_D$ is realised by a
feasible design, and if it is not a constant then no point of $S$ is feasible. Consequently the per-chart
test of Algorithm~\ref{alg:decide} is exact, and the algorithm outputs exactly \textbf{feasible} or
\textbf{infeasible}.
\end{claim}

\begin{proof}
Claim~\ref{clm:admissible} supplies \textbf{(H3)} of Lemma~\ref{lem:genericgcd} with $d_0=\omega$:
after clearing, the $\lambda$-leading coefficient of $m_{I,J}$ is $\det H_{I,J}$ times the clearing factor,
hence a unit of $K[s,1/\delta]$, and it is a non-zero constant for the uncleared pencil; it also makes
every point of $S$ admissible. Hypotheses \textbf{(H1)}, \textbf{(H2)} and \textbf{(H4)} of the lemma hold because a
chart of $\mathcal I_D$ has as its parameter space the affine subspace of $\mathbb R^{j(\dim V_\varphi-j)}$
cut out by the linear condition $\Phi_0\subseteq\Phi$ (Definition~\ref{defn:charts}), and $S$ is its
non-empty Zariski-open chart proper, so a non-zero polynomial cannot vanish on all of $S$, which is
\textbf{(H4)}.
Clause~(2) of the lemma is the
negative direction: a non-constant generic gcd makes every specialised gcd non-constant, so every pencil in
the chart is singular at some $\lambda$. For the positive direction, clause~(3) of the lemma produces a real
point $s\in S$ at which the specialised gcd is $1$; by Theorem~\ref{thm:reduction}(4) this gives
$\dim(R_\lambda+\Phi(s))=\omega$ for every $\lambda$, and by Theorem~\ref{thm:realize} the subspace
$\Phi(s)\in\mathcal I_D$ is realised by some $W$, so the design $V=D^{\perp_{K_0}}\oplus W$ is feasible. The
last sentence is the definition of the algorithm: it returns \emph{feasible} at the first chart whose generic
gcd is constant, and \emph{infeasible} only when every chart fails.
\end{proof}

\begin{thm}[correctness]\label{thm:correct}
Assume (i) $c=2$, (ii) simple spectrum over $\mathbb C$ of $M$, and (iii) the split condition. Then
Algorithm~\ref{alg:decide} terminates,
and it returns \emph{feasible} if and only if a feasible $K_0$-split design of order $q$ with prescribed rows
$V_0$ exists.
\end{thm}

\begin{proof}
Write \textbf{Feas} for the assertion that a feasible $K_0$-split design of order $q$ with prescribed rows
$V_0$ exists. We prove the three equivalences (a)--(c) below; composing them gives the theorem, and in each
of them the support is named at the start.

\emph{(a) \textbf{Feas} holds if and only if there are $D\in\mathcal B_{\rm pres}$ and a subspace $W$ with
$W_0\subseteq W\subseteq T_D$, $\dim W=t$, and}
\begin{equation}\label{eq:chainA}
R_\lambda+\Phi(W)=D^\circ\quad\text{for all }\lambda .
\end{equation}
\emph{Support: Theorem~\ref{thm:splitfibre} (the fixed-kernel fibre) and Theorem~\ref{thm:reduction}
(image-space reduction).} If $\mathcal L$ is a feasible split design and $D=\ker H$, then $D\in\mathcal I_c(M)$
by Claim~\ref{clm:invariant}, and the condition $D\in\mathcal B_{\rm pres}$ is exactly (H1) of
\eqref{eq:splithyp}; Theorem~\ref{thm:splitfibre} exhibits the row space of $\mathcal L$ as
$V=D^{\perp_{K_0}}\oplus W$ with $W_0\subseteq W\subseteq T_D$ and $\dim W=t$. Conversely, given such $D$ and
$W$, the row space $V=D^{\perp_{K_0}}\oplus W$ has dimension $(k-c)+t=\omega+t=q$, contains $V_0$ by (H1)
and (H2), is split by construction, and satisfies \eqref{eq:5a} by Theorem~\ref{thm:splitfibre}. In both
directions condition \eqref{eq:5b} is equivalent to \eqref{eq:chainA}: this is part~(4) of
Theorem~\ref{thm:reduction}, whose row-space identity reads
$\row(\lambda H-H_A)=R_\lambda+\Phi(W)$.

\emph{(b) Condition (a) holds if and only if there are $D\in\mathcal B_{\rm pres}$ and $\Phi\in\mathcal I_D$
with}
\begin{equation}\label{eq:chainB}
\dim\bigl(R_\lambda+\Phi\bigr)=\omega\quad\text{for all }\lambda .
\end{equation}
\emph{Support: Theorem~\ref{thm:realize} (realisability) and Theorem~\ref{thm:reduced} (the reduced
criterion).} By Theorem~\ref{thm:realize}, a subspace $\Phi\le D^\circ$ equals $\Phi(W)$ for some $W$ with
$W_0\subseteq W\subseteq T_D$ and $\dim W=t$ if and only if $\Phi\in\mathcal I_D$, so the quantifier
$\exists W$ of (a) may be replaced by $\exists\Phi\in\mathcal I_D$; and since
$R_\lambda+\Phi\subseteq D^\circ$ and $\dim D^\circ=k-c=\omega$, condition \eqref{eq:chainA} on the
subspace $\Phi$ is exactly \eqref{eq:chainB}. This is the content of Theorem~\ref{thm:reduced}, whose second
form \eqref{eq:reduced2} is not needed here.

\emph{(c) Condition (b) holds if and only if Algorithm~\ref{alg:decide} returns \emph{feasible}.}
\emph{Support: Claim~\ref{clm:admissible} (admissibility), Lemma~\ref{lem:genericgcd} (generic gcd) and
Theorem~\ref{thm:base} (classification of the base).} By Theorem~\ref{thm:base} the set $\mathcal I_2(M)$ is
finite and enumerated explicitly in step~2, so only finitely many $D$ are tried and
$\mathcal B_{\rm pres}$ is recognised by the linear condition \eqref{eq:bpres}; by
Definition~\ref{defn:charts} the set $\mathcal I_D$ is covered by finitely many charts and every
$\Phi\in\mathcal I_D$ lies on one of them. Thus condition (b) is equivalent to the assertion that some chart
$S$ of $\mathcal I_D$ contains a real point $s$ for which \eqref{eq:chainB} holds at $\Phi(s)$.
On a chart $S$, condition \eqref{eq:chainB} at $\Phi(s)$ says that the maximal minors of the pencil at $s$
have no common zero; every point of $S$ is admissible by Claim~\ref{clm:admissible}, so
Lemma~\ref{lem:genericgcd} applies: the generic gcd of $S$ is a non-zero constant exactly when such a real
point exists, and a non-constant generic gcd rules out every point of $S$. The gcd test of steps~5--6
therefore returns \emph{feasible} exactly when (b) holds, and \emph{infeasible} exactly when it fails.
Finitely many $D$ and finitely many charts per $D$ give termination.
\end{proof}

Correctness being established, the only question left is the cost of the enumeration and of the gcd
computations, which is Part D.

\subsection{Bit complexity}\label{sec:complexity}

\emph{Input length.} We recall from the Notation that $\tau$ is the maximum bit length of an entry of $A$,
$C$ or $L_0$, so that the input is encoded by $O(n^2\tau)$ bits and every estimate below is a function of
$n$ and $\tau$ alone. The prescribed block is part of the input length and not an afterthought: the vectors
of $V_0=\row L_0$ enter $W_0=P_{K_0^\perp}(V_0)$, $w_0=\dim W_0$, $\Phi_0=\varphi(W_0)$, the range
$w_0\le t\le(n-k)-r(D)$ of the fibre parameter, and the membership test $D\perp P_{K_0}(V_0)$ that defines
$\mathcal B_{\rm pres}$ in \eqref{eq:bpres}. A bound in the bit length of $A$ and $C$ alone would therefore
bound neither the enumeration of the base nor the gcd test per chart, and no such bound is claimed anywhere
below.

\emph{Field arithmetic.} In branch~R the coefficients lie in $L_D=\mathbb Q(\alpha,\beta)$, in branch~C in
$L_D=\mathbb Q(\mu+\bar\mu,\mu\bar\mu)$; in both branches $d_D:=[L_D:\mathbb Q]\le k^2$
(Theorem~\ref{thm:base}). We use the standard exact representation of \cite{cohen93}: $L_D=\mathbb Q[t]/(F)$, with $F$
the minimal polynomial of a primitive element, an element being given by
its coordinate vector of length $d_D$ over $\mathbb Q$. In this representation addition costs $O(d_D)$
coefficient operations, and multiplication and inversion cost $O(d_D^2)$ (schoolbook; $\widetilde O(d_D)$
with fast polynomial arithmetic), each coefficient operation being integer arithmetic on numbers of bit
length $\poly(n,\tau)$. The cost therefore depends on three input-size parameters: the degree $d_D$, the
height $h(F)$ of the defining polynomial, and the height of the coordinates; all three are $\poly(n,\tau)$
by Appendix~\ref{app:height}, Steps~$8$--$8'$. A bound on $d_D$ alone --- such as $d_D\le k^2$ --- controls
neither $h(F)$ nor the cost of one field operation, and by itself would give no polynomial time bound.

\emph{One determinant bound.} All height propagation below rests on a single inequality, proved in
Appendix~\ref{app:height}, Step~9: if $B$ is an $r\times r$ matrix whose entries have height at most $H$,
then $h(\det B)\le rH+\tfrac{r}{2}\log_2 r+O(r)$; the same bound applies to the cofactors out of which
kernels, orthogonal complements and subresultants are built.

\begin{claim}[arithmetic of the enumeration]\label{clm:arith}
Assume simple spectrum over $\mathbb C$. Then:
\begin{enumerate}[label=(\roman*),leftmargin=*,itemsep=1pt,topsep=2pt]
\item the number of planes is $O(k^2)$, and they are enumerated in $\poly(n,\tau)$ bit operations, each
plane costing polynomially many field operations on data of height $\poly(n,\tau)$;
\item each plane is defined over a number field $L_D$ of degree $d_D:=[L_D:\mathbb Q]\le k^2$
(branch~R: $L_D=\mathbb Q(\alpha,\beta)$; branch~C: $L_D=\mathbb Q(\mu+\bar\mu,\mu\bar\mu)$),
represented as $\mathbb Q[t]/(F)$ with $h(F)=\poly(n,\tau)$; the
$\mathbb Q$-basis $1,t,\dots,t^{d_D-1}$ consists of elements of length $\poly(n,\tau)$, and addition,
multiplication and inversion in $L_D$ cost $O(d_D)$, $O(d_D^2)$ and $O(d_D^2)$ coefficient operations
respectively, each on integers of height $\poly(n,\tau)$;
\item the field data and the bases of $D$, $D^{\perp_{K_0}}$, $T_D$, together with $V_0=\row L_0$,
$W_0=P_{K_0^\perp}(V_0)$, $w_0$, $\varphi$, $\Phi_0$, the membership test defining $\mathcal B_{\rm pres}$ and the
entries of the cleared pencil, are computed in $\poly(n,\tau)$ bit operations and have heights
$\poly(n,\tau)$.
\end{enumerate}
Here $\tau$ is the bit length of the whole input $(A,C,L_0)$.
\end{claim}

\begin{proof}[Sketch]
$M$ is rational, so its characteristic polynomial has coefficients of polynomial bit length; factoring it
over $\mathbb Q$ is polynomial-time in the bit size of the input by the LLL-based algorithm of \cite{lll82}
(see also \cite{cohen93}), and its irreducible factors have polynomial height by Mignotte's bound; the
roots have height bounded by the Mahler measure of these factors. Branch R uses $L_D=\mathbb Q(\alpha,\beta)$ with
$[\mathbb Q(\alpha,\beta):\mathbb Q]\le d_id_j\le k^2$; a primitive element $\theta=\alpha+\kappa\beta$ with integral
$\kappa\le d_id_j$ has a minimal polynomial $F$ of degree $d_D\le k^2$ and height $\poly(n,\tau)$ by the
resultant bound of Step~$8'$; eigenvectors are obtained by Cramer's rule over
$L_D$, with denominators $f_i'(\alpha)\ne0$ by simple spectrum over $\mathbb C$. Branch C works over
$L_D=\mathbb Q(\mu+\bar\mu,\mu\bar\mu)\subseteq\mathbb Q(\mu,\bar\mu)$, of degree $d_D\le d^2\le k^2$
(Theorem~\ref{thm:base}); its field data are the two algebraic numbers $\mu+\bar\mu$ and $\mu\bar\mu$,
both of height $\poly(n,\tau)$ by Step~7, and Step~$8'$ supplies $F$ with $h(F)=\poly(n,\tau)$ from
them, so the final objects live over $L_D$, not over $\mathbb Q$ in general. The kernel itself is
computed from an
integral model of $[C;CA]$ (Step~2), so that $N$ is a primitive integral matrix of full column rank $k$ and
$P_{K_0}=N(N^{\mathsf T}N)^{-1}N^{\mathsf T}$ is the orthogonal projection onto $K_0$; orthogonal
complements and kernels are again Cramer, the prescribed data $W_0$, $w_0$, $\Phi_0$ and the membership test of
$\mathcal B_{\rm pres}$ are computed from $L_0$ as in Step~$4'$, and the pencil entries are obtained by
polynomially many field operations.
\end{proof}

The height propagation is checked step by step in Appendix~\ref{app:height}.

\begin{claim}[degrees and size on an image chart]\label{clm:degrees}
On a chart of $\Gr(j,V_\varphi)$ at most $j$ rows of the pencil depend on the chart parameters, and each
dependence is affine. Hence every maximal minor satisfies
\[
\deg_{\text{param}}\le j\le\min(t,\,2p),\qquad \deg_\lambda\le\omega\le k ,
\]
and, after clearing the chart denominator, its coefficients have height at most $H=\poly(n,\tau)$, with
$O((j+1)^{p^2})=\poly(n,\tau)$ coefficients per minor. The chart carries
\[
\nu:=\binom{q}{\omega}\binom{k}{\omega}
=\binom{\omega+t}{t}\binom{k}{2}=O\bigl(k^{\,t+2}\bigr)\le O\bigl(k^{\,2p+2}\bigr)
\]
maximal minors, because $t\le\dim T_D=(n-k)-r(D)\le n-k\le2p$ on a nonempty chart. In particular on a
line chart ($j=1$) the parameter degree is at most one, and once $p$ is fixed the chart carries
polynomially many minors of bounded parameter degree.
\end{claim}

\begin{proof}
A maximal minor is multilinear in the rows; exactly the $j$ rows corresponding to a basis of $\Phi$ vary,
and they vary affinely in the chart coordinates, so the total parameter degree is at most $j$. The row and
column index sets range over $\binom{q}{\omega}$ and, since $\omega=k-2$ at $c=2$, over $\binom{k}{2}$,
which gives $\nu$; the bound $t\le2p$ is the non-emptiness condition $t\le\dim T_D$ of
Theorem~\ref{thm:splitfibre} combined with $\dim T_D\le n-k\le2p$. Heights are those of
Claim~\ref{clm:arith} and of Appendix~\ref{app:height}, Step~9.
\end{proof}

\begin{thm}[bit complexity]\label{thm:complexity}
Assume (i) $c=2$, (ii) simple spectrum over $\mathbb C$, (iii) the split condition, and (iv) \textbf{the
output dimension $p$ is fixed}, so that $n-k=\rank[C;CA]\le2p$ is a constant. Then
Algorithm~\ref{alg:decide} runs in
$\poly(n,\tau)$ bit operations, where $\tau$ is the bit length of the whole input $(A,C,L_0)$
(Notation and the beginning of this subsection), and all arithmetic takes place in number fields of degree
at most $k^2$. The bound is coarse: the exponents are not optimized and the factor $\poly(n,\tau)$ hides
constants that depend on the fixed number $p$.
\end{thm}

\begin{proof}
The cost has six factors, one per stage of the algorithm. Each is bounded separately and the product is the
stated coarse bound; no attempt is made to exploit that the stages share data.
\begin{enumerate}[label=(\alph*),leftmargin=*,itemsep=1pt,topsep=2pt]
\item \emph{The base.} By Theorem~\ref{thm:base} and Claim~\ref{clm:arith}(i) there are $O(k^2)$ planes,
and their enumeration --- factorization of $\chi_M$ included --- costs $\poly(n,\tau)$ bit operations, so
this factor is $O(k^2)$ planes at $\poly(n,\tau)$ each.
\item \emph{The prescribed rows.} Membership in $\mathcal B_{\rm pres}$ and the data $W_0$, $w_0$,
$\Phi_0$, $T_D$ are computed from $L_0$ once per plane, at cost $\poly(n,\tau)$ each
(Claim~\ref{clm:arith}(iii) and Appendix~\ref{app:height}, Step~$4'$).
\item \emph{The charts.} For each plane, $j$ ranges over at most $\min(t,2p)+1$ values, and the standard
affine charts of $\Gr(j,V_\varphi)$ number at most $\sum_j\binom{\dim V_\varphi}{j}\le2^{2p}$; each chart
carries $j(\dim V_\varphi-j)\le p^2$ parameters (Remark~\ref{rem:imagesize}). With $p$ fixed this is a
constant number of charts of constant dimension.
\item \emph{The minors.} By Claim~\ref{clm:degrees} a chart carries $\nu=O(k^{2p+2})$ maximal minors, each
of $\lambda$-degree at most $\omega\le k$, of parameter degree at most $j\le2p$, with coefficient height at
most $H=\poly(n,\tau)$ (Claim~\ref{clm:arith} and Appendix~\ref{app:height}, Step~9).
\item \emph{The gcd chain.} The gcd of the minors over $K(s)[\lambda]$ is computed by the subresultant
algorithm \cite{vzgg13}: two $\lambda$-polynomials of degree at most $\omega$ are reduced in $O(\omega^2)$
ring operations, and every intermediate subresultant has bit length at most
$\varrho=O(\omega(H+\log\omega))$ by the Hadamard bound of Appendix~\ref{app:height}, Step~10. Iterating
over the $\nu$ minors therefore costs $O(\nu\omega^2)$ ring operations, each on integers of bit length
$\varrho$, and each ring operation in a field coefficient costs a further factor $O(d_D^2)$ with
$d_D=[L_D:\mathbb Q]\le k^2$ (\cite{cohen93}; Appendix~\ref{app:height}, Step~$8'$).
\item \emph{Summation.} Multiplying the factors, the number of bit operations is at most
\[
\begin{gathered}
\underbrace{O(k^2)}_{\text{planes}}\;
\underbrace{\poly(n,\tau)}_{\text{per-plane data}}\;
\underbrace{O(2^{2p}(p+1))}_{\text{charts}}\;
\underbrace{O(k^{2p+2})}_{\text{minors}}\\[2pt]
{}\times{}\;
\underbrace{O(\omega^{2})}_{\text{gcd steps}}\;
\underbrace{O(d_D^{2})}_{\text{field ops}}\;
\underbrace{O(\varrho^{2})}_{\text{bits per ring op}} ,
\end{gathered}
\]
which is $\poly(n,\tau)$ because $p$ is fixed; summing over the $O(k^2)$ planes is already included in the
first factor.
\end{enumerate}
\end{proof}

\begin{remark}[the two boundaries, and why $p$ must be fixed]\label{rem:boundaries}
Hypothesis (iv) of Theorem~\ref{thm:complexity} is not cosmetic. Without a bound on
$n-k=\rank[C;CA]$ the space $V_\varphi=\varphi(T_D)$ grows with $k$, the fibre parameter has up to
$p^2$-many free parameters, and the chartwise gcd is no longer a polynomial-time primitive. Note also
that $c=2$ alone gives $\omega=k-2$, so $D^\circ$ is \emph{large}: it is the bound on $n-k$, not the
smallness of $c$, that controls the fibre. Likewise, without simple spectrum over $\mathbb C$ the base is
positive-dimensional (a Grassmannian of dimension $2(d_\lambda-2)=O(n)$) and the number of charts is no
longer polynomial. With $c$ allowed to vary, the number of components is $\binom{k}{c}$, which is
exponential. The present method therefore gives $\poly(n,\tau)$ only for the split problem, under fixed
$p$ and simple spectrum over $\mathbb C$ --- with $\tau$ the bit length of the whole input $(A,C,L_0)$, as
in the Notation --- and at best $n^{O(cn+p^2)}$ in general.
\end{remark}

\subsection{A worked example}\label{sec:example}

This subsection carries one instance through the reduction and the certificate: the compressed data,
the admissible kernel direction, the image space, the pencil, and the gcd certificate. All entries are
exact integers, and the computation is the one behind the $c=2$ side of Theorem~\ref{thm:nonmono},
with the intermediate matrices displayed.

\emph{The instance.} $n=8$, $p=2$, and
\begin{equation}\label{eq:exA}
\begin{gathered}
A=\begin{pmatrix}
0 & 0 & 1 & 0 & 0 & 0 & 0 & 0\\
0 & 0 & 0 & 1 & 0 & 0 & 0 & 0\\
1 & -2 & 1 & 0 & -1 & 1 & 1 & -1\\
-1 & 2 & 1 & 0 & 2 & -2 & -1 & 2\\
0 & 0 & 0 & 0 & 1 & 0 & 0 & 0\\
0 & 0 & 0 & 0 & 0 & 1 & 0 & 0\\
0 & 0 & 0 & 0 & 0 & 0 & 2 & 0\\
0 & 0 & 0 & 0 & 0 & 0 & 0 & 3
\end{pmatrix},\\
C=\begin{pmatrix}
1&0&0&0&0&0&0&0\\ 0&1&0&0&0&0&0&0
\end{pmatrix}.
\end{gathered}
\end{equation}
Here $\rank[C;CA]=4$, so $k=4$, and
\[
K_0=\ker[C;CA]=\operatorname{span}\{e_5,e_6,e_7,e_8\},\qquad
N=\begin{pmatrix}0_{4\times4}\\ I_4\end{pmatrix}.
\]
The compression is $M=\pi_{K_0}A|_{K_0}=\operatorname{diag}(1,1,2,3)$ and the coupling is
\[
A_{21}=\begin{pmatrix}
0&0&0&0\\ 0&0&0&0\\ -1&1&1&-1\\ 2&-2&-1&2\\ 0&0&0&0\\ 0&0&0&0\\ 0&0&0&0\\ 0&0&0&0
\end{pmatrix}.
\]
The coupling has rank two, but its restriction to the direction below has rank one.

\emph{The direction.} In the coordinates of $K_0$ fixed by $N$, take
$D=\operatorname{span}\{e_1+e_2,\,e_3\}\le K_0$. Then $M(e_1+e_2)=e_1+e_2$ and $Me_3=2e_3$, so
$M(D)\subseteq D$ and $D\in\mathcal I_2(M)$ with $\dim D=c=2$. Since $L_0=0$ the prescribed-row
conditions read $P_{K_0}(V_0)=0$ and $W_0=0$, which hold trivially, so $D\in\mathcal B_{\rm pres}$.

\emph{The fibre.} Here $A_{21}D=\operatorname{span}\{(0,0,1,-1,0,0,0,0)^{\mathsf T}\}$, so
\[
r(D)=1,\qquad \dim T_D=(n-k)-r(D)=4-1=3.
\]
With $t=1$ (i.e.\ $q=\omega+t=2+1=3$), $w:=(0,0,1,1,0,0,0,0)^{\mathsf T}$ and
$W=\operatorname{span}\{w\}\subseteq T_D$, one computes
\[
\varphi_w=(1,-1,0,1)\in D^\circ,\qquad \Phi(W)=\operatorname{span}\{(1,-1,0,1)\},
\]
so $\dim\Phi(W)=1=\min(t,2p)$: the image parameter is a line, as in Definition~\ref{defn:charts}. The
resulting design row space is
\[
\begin{aligned}
V=D^{\perp_{K_0}}\oplus W
&=\operatorname{span}\bigl\{(0,0,0,0,-1,1,0,0),\\
&\qquad (0,0,0,0,0,0,0,1),\\
&\qquad (0,0,1,1,0,0,0,0)\bigr\},
\end{aligned}
\]
of dimension $q=3$; the design is $K_0$-split by construction.

\emph{The pencil.} With these three vectors as the rows of $\mathcal L$,
\[
\begin{gathered}
H=\mathcal L N=\begin{pmatrix}-1&1&0&0\\ 0&0&0&1\\ 0&0&0&0\end{pmatrix},\\
H_A=\mathcal L AN=\begin{pmatrix}-1&1&0&0\\ 0&0&0&3\\ 1&-1&0&1\end{pmatrix}.
\end{gathered}
\]
The invariants are $\omega=\rank H=2$, $\dim(K_0\cap V^\perp)=2=c$ (exact kernel), and
$\row H_A\subseteq\row H$, i.e.\ \eqref{eq:5a} holds. The nonzero maximal $2\times2$ minors of
$\lambda H-H_A$ are
\[
\begin{gathered}
-(\lambda-3)(\lambda-1),\quad (\lambda-3)(\lambda-1),\quad \lambda-1,\\
1-\lambda,\quad \lambda-3,\quad 3-\lambda .
\end{gathered}
\]

\emph{The certificate.} The minors $\lambda-1$ and $\lambda-3$ are coprime, with Bézout identity
\[
\tfrac12\bigl((\lambda-1)-(\lambda-3)\bigr)=1 ,
\]
so the gcd of all maximal minors is $1$ and condition~\eqref{eq:5b} holds. This is a direct
certificate for the displayed design; it does not require the algorithm. Indeed a feasible design of
order $q=3$ with $c=2$ exists, whereas by Theorem~\ref{thm:nonmono} no feasible design of order $3$
with $c=1$ does.

\emph{Scope.} The eigenvalue $1$ of $M$ has a two-dimensional eigenspace, so $\mathcal I_2(M)$ is
positive-dimensional (Remark~\ref{rem:nonsimple}) and the finite enumeration of
Theorem~\ref{thm:base} does not apply. The proofs of Theorems~\ref{thm:correct} and
\ref{thm:complexity} use simple spectrum over $\mathbb C$, i.e.\ distinct eigenvalues in $\mathbb C$;
this instance therefore illustrates the reduction and the certificate format, not the algorithm's
execution or those proofs. The distinction is worth recording, because it shows that feasibility itself is
not confined to the simple-spectrum case --- only the correctness and cost arguments are.

\section{Discussion and conclusion}\label{sec:conclusion}

This section collects what the paper does not claim and records the evidence behind what it does claim.
Part A consolidates the limitations, Part B the evidence.

\subsection{Limitations}\label{sec:conclusion:limits}

\begin{enumerate}[label=\arabic*.,leftmargin=*,itemsep=1pt,topsep=2pt]
\item All structural results are for the split problem \textbf{S}; the unrestricted problem \textbf{U} is
only reduced to a fibre description (Theorem~\ref{thm:ufibre}) and its invariant-kernel structure is
lost.
\item The decision algorithm and the complexity bound require simple spectrum over $\mathbb C$; the
non-simple case is left open, and the worked example of Section~\ref{sec:example}, whose compression has
a repeated eigenvalue, illustrates the certificate format rather than these two results.
\item No uniformity in $c$, and none in $p$, is claimed; the honest general bound is $n^{O(cn+p^2)}$. The
polynomial bound of Theorem~\ref{thm:complexity} is $\poly(n,\tau)$ with the output dimension $p$ fixed and
with $\tau$ the bit length of the whole input $(A,C,L_0)$; in particular it is \emph{not} a bound in the bit
length of $A$ and $C$ alone, which would omit the data $V_0$, $W_0$, $\Phi_0$ and $\mathcal B_{\rm pres}$
that the algorithm reads off $L_0$.
\item The general invariant-plane variety is used only through its defining equations; no rational
parametrization is claimed.
\end{enumerate}

\emph{The promise, precisely.} What the paper establishes is this: for the $K_0$-split problem at $c=2$,
if $M$ has simple spectrum over $\mathbb C$ and the output dimension $p$ is fixed, then the exact yes/no
decision algorithm of Section~\ref{sec:decision} terminates and runs in $\poly(n,\tau)$ bit operations,
with $\tau$ the bit length of the whole input $(A,C,L_0)$ (Theorems~\ref{thm:correct}
and~\ref{thm:complexity}). Each hypothesis is used: the split condition is what makes $D$ an invariant
plane (Claim~\ref{clm:invariant}), simple spectrum over $\mathbb C$ is what makes the base finite
(Theorem~\ref{thm:base}), and a fixed $p$ is what keeps the dimension of the image parameter bounded
(Remark~\ref{rem:boundaries}). No claim is made for the unrestricted problem at $c=2$, nor for the split
problem at $c=2$ when the spectrum is not simple over $\mathbb C$, nor for $c\ge3$.

The disclaimers inside the body are part of the same statement and are not weakened here:
Remark~\ref{rem:nonsimple} is the disclaimer attached to the classification of
Theorem~\ref{thm:base}, Remark~\ref{rem:boundaries} lists the two boundaries of the complexity bound, and
Remark~\ref{rem:open} lists the two questions on general $c$ that we do not answer: (a) whether the
quantifier over $\Phi$ can be replaced by a rank or separation condition on the induced maps of
$T_D/W_0$, and (b) whether the layer index can be traded for $\dim U$ with $U=\pi_1(W)$, since
$c=k-\dim U$. Nothing is claimed about the complexity of the general decision problem.

\subsection{Reproducibility and evidence}\label{sec:conclusion:evidence}

This paper has \emph{no} machine-checked formalisation. Unlike companion manuscripts that carry Lean
developments, no statement of Sections~\ref{sec:layers}--\ref{sec:decision} is formalised at the time of
writing, and nothing in those sections should be read as machine-verified.

The evidence behind the concrete assertions is of two kinds. First, the proofs of
Sections~\ref{sec:layers}--\ref{sec:decision} are paper-level and complete, with one exception:
the height chain of Claim~\ref{clm:arith} is a sketch in the body, made explicit step by step in
Appendix~\ref{app:height}; the constants there are standard and deliberately not optimized, so no
single numerical constant is claimed. Two constants are not pinned down at all, and are marked in the
appendix as coarse bounds rather than asserted: the multiplicative constant of the subresultant chain
(Step~10) and the exponent hidden in the factor $\poly(n,\tau)$ of the enumeration
(Claim~\ref{clm:arith}). The bit length $\tau$ entering the bound is that of the whole input
$(A,C,L_0)$, the prescribed block included. Second, the three places where the paper asserts something
about a specific system --- the instance of Theorem~\ref{thm:nonmono}, the worked example of
Section~\ref{sec:example}, and the classification of Table~\ref{tab:planes} --- rest on exact
rational arithmetic, not on floating-point computation. In the proof of Theorem~\ref{thm:nonmono}
the nonzero maximal $2\times2$ minors are expanded over $\mathbb Q$ and their gcd is $1$ precisely
because the displayed combination $\tfrac12\bigl((\lambda-1)-(\lambda-3)\bigr)$ is a rational
combination of two of them, and on the $c=1$ side the gcd $(\lambda-1)(\lambda-2)(\lambda-3)$ is
read off the displayed integer matrix $H$ for symbolic $s$; both computations involve no algebraic
extension and no rounding. In the worked example of Section~\ref{sec:example} the data $A,C$ are
integral, and $N$, $M$, $A_{21}$, the design basis, $H$, $H_A$, all $2\times2$ minors of the pencil
and their gcd are computed over $\mathbb Q$; six minors are nonzero, the displayed Bézout identity
is exact, and again no algebraic extension or rounding occurs, so the certificate is reproducible in
any exact arithmetic system. Table~\ref{tab:planes} is a classification rather than a computation:
its two counts are the combinatorial expressions of Theorem~\ref{thm:base}, and its third column is
the field statement of Claim~\ref{clm:arith}.

\appendices
\small
\section{Proof of the height propagation}\label{app:height}

This appendix records the chain used in Claim~\ref{clm:arith} with the standard bounds made
explicit. Throughout, $h(\cdot)$ is the maximum bit length of the entries of a matrix or of the
coefficients of a polynomial, and $\tau$ is the bit length of the whole input $(A,C,L_0)$ as in the
Notation.

\emph{Step 0 (input).} $h(A),h(C),h(L_0)\le\tau$.

\emph{Step 1 ($CA$ and $[C;CA]$).} Entrywise, $(CA)_{ij}$ is a sum of $n$ products, so
$h(CA)\le 2\tau+\lceil\log_2 n\rceil$. Hadamard's bound gives, for every square submatrix $S$ of
$[C;CA]$ of order $d$, $|\det S|\le(\sqrt{n}\,2^{2\tau+\log n})^{d}$, i.e.\
$h(\det S)\le d\,(2\tau+\log n+\tfrac12\log n)$. Only $A$ and $C$ enter this step; the prescribed block
enters in Step~$4'$.

\emph{Step 2 (integerisation and the kernel basis $N$).} A kernel does not change under multiplication by a
non-zero scalar, so we first clear denominators: let $\Delta_0$ be a common denominator of the entries of
$[C;CA]$, so that $\log_2\Delta_0\le(2p+n)\,(2\tau+\lceil\log_2 n\rceil)$ and
$M_0:=\Delta_0[C;CA]$ is integral with $\ker M_0=\ker[C;CA]$. A basis $N$ of $\ker M_0$ can be chosen so
that each entry is a cofactor of $M_0$, hence of order at most $2p$; dividing $N$ by the gcd of its entries
(an integer of no larger height) makes $N$ primitive and changes no kernel. By Step~1 and the Hadamard
bound for a cofactor of order $2p$,
\[
h(N)\le 2p\,\bigl(h([C;CA])+h(\Delta_0)+\tfrac12\log_2(2p)\bigr)+O(p),
\]
which is $\poly(n,\tau)$. Since $N$ is a basis of $K_0$, it has full column rank $k$; hence
$N^{\mathsf T}N$ is positive definite and invertible, and $P_{K_0}=N(N^{\mathsf T}N)^{-1}N^{\mathsf T}$ is
the orthogonal projection onto $K_0$. This full-column-rank property is what the formula requires, and it
is used in Steps~3--4; replacing $N$ by $NC$ with $C$ invertible leaves $P_{K_0}$ unchanged.

\emph{Step 3 (Gram matrix and its inverse).} $G=N^{\mathsf T}N$ has entries that are sums of $n$
products, so $h(G)\le 2h(N)+\log n$. Since $G$ is a positive definite integer matrix, $\det G\ge1$,
and the entries of $G^{-1}$ are cofactors over $\det G$, so
\[
h(G^{-1})\le k\,(2h(N)+\log n+\tfrac12\log k)+O(k).
\]

\emph{Step 4 (compression and coupling).} $M=(N^{\mathsf T}N)^{-1}N^{\mathsf T}AN$ and
$A_{21}=P_{K_0^\perp}AN$ are obtained from $G^{-1}$, $N$ and $A$ by polynomially many ring
operations, so
\[
h(M),\,h(A_{21})\le 4h(N)+2h(G^{-1})+\tau+O(\log n),
\]
which is $\poly(n,\tau)$.

\emph{Step 4$'$ (the prescribed block and the data derived from it).} The vectors of $V_0=\row L_0$ have
entries of height at most $\tau$ by Step~0. The projection
$W_0=P_{K_0^\perp}(V_0)$ is the row space of $L_0\bigl(I-N(N^{\mathsf T}N)^{-1}N^{\mathsf T}\bigr)$ and is obtained from $L_0$, $N$
and $G^{-1}$ by $O(n^3)$ ring operations, so
\[
h(W_0)\le h(L_0)+2h(N)+2h(G^{-1})+O(\log n)=\poly(n,\tau),
\]
and the same bound holds for $w_0=\dim W_0$, for the matrix of $\varphi$ (which is $A_{21}$ read through
the identification $\iota_N$ of \eqref{eq:iotaN}) and for $\Phi_0=\varphi(W_0)$. The condition
$D\perp P_{K_0}(V_0)$ defining $\mathcal B_{\rm pres}$ \eqref{eq:bpres} is a rank test on a matrix whose
entries are inner products of the rows of $L_0$ with those of $N$; it therefore adds no height beyond
$h(L_0)+h(N)$ and costs $\poly(n,\tau)$ bit operations per plane.

\emph{Step 5 (characteristic and minimal polynomial).} The coefficient of $t^{\,k-i}$ in $\chi_M$ is
a sum of principal $i\times i$ minors, hence at most $\binom{k}{i}(\sqrt{k}\,2^{h(M)})^{i}$ in
absolute value; so $h(\chi_M)\le k\,h(M)+O(k\log k)$. The minimal polynomial divides $\chi_M$.

\emph{Step 6 (irreducible factors).} By Mignotte's bound, a divisor $f$ of a degree-$k$ integer
polynomial $g$ satisfies $\|f\|_\infty\le\binom{k}{\lfloor k/2\rfloor}\|g\|_\infty$, so
\[
h(f)\le h(\chi_M)+k .
\]

\emph{Step 7 (roots).} Represent a root $\alpha$ of $f$ by the class of $t$ in
$\mathbb Q[t]/(f)$; its height is bounded in terms of $h(f)$. Cauchy's bound gives
$|\alpha|\le2\|f\|_\infty^{1/\deg f}$, and the logarithmic height satisfies
$h(\alpha)\le h(f)+\log\deg f$.

\emph{Step 8 (eigenvectors).} Solving $(M-\alpha I)v=0$ over $L=\mathbb Q(\alpha)$ by Cramer
expresses each coordinate of $v$ as a quotient of two minors of order $k-1$ of $M-\alpha I$. Those
minors are polynomials in $\alpha$ of degree at most $k-1$ whose coefficients involve $k!$ products
of entries of $M$; hence
\[
h(v)\le (k-1)\bigl(h(M)+h(\alpha)\bigr)+O(k\log k).
\]
The denominators are nonzero because the eigenvalues are distinct.

\emph{Step 8$'$ (the field $L_D$ and the cost of its arithmetic).} In branch~R
the field is generated by a primitive element $\theta=\alpha+\kappa\beta$ with integral $\kappa\le d_id_j\le k^2$,
whose minimal polynomial $F$ may be taken to be a factor of
$\operatorname{Res}_y\bigl(f_i(y),f_j(\theta-y)\bigr)$; the resultant is a determinant of a Sylvester
matrix of order $d_i+d_j\le2k$ whose entries have height at most $h(f_i)+h(f_j)$, so Hadamard's bound
gives it height $O\bigl(k\,(h(f_i)+h(f_j)+\log k)\bigr)$, and $F$ has no larger height by Mignotte's bound
(Step~6). Hence in branch~R $d_D=\deg F\le k^2$ and $h(F)=\poly(n,\tau)$. In branch~C the field is
$L_D=\mathbb Q(\mu+\bar\mu,\mu\bar\mu)\subseteq\mathbb Q(\mu,\bar\mu)$ (Theorem~\ref{thm:base}); its two
generators have height $\poly(n,\tau)$ by Step~7 and the standard inequalities for sums and products of
algebraic numbers, and an algebraic number of degree at most $k$ and
height $H$ has a minimal polynomial of height at most $k\,(H+1)+O(k)$, so the same primitive-element
construction, applied to their two minimal polynomials, gives $F$ of degree $d_D\le d^2\le k^2$ and
height $\poly(n,\tau)$. Thus $d_D\le k^2$ and $h(F)=\poly(n,\tau)$ in either branch.
Elements of $L_D$ are coordinate vectors of
length $d_D$ in the $\mathbb Q$-basis $1,\theta,\dots,\theta^{d_D-1}$; by Cauchy's bound
$|\theta|\le1+\|F\|_\infty$, so the $i$-th basis element has length $(1+\|F\|_\infty)^{i}$, and a field
element whose coordinates have height $H_0$ has length at most
$d_D\,H_0\,(1+\|F\|_\infty)^{d_D}$, both $\poly(n,\tau)$. In this representation addition costs $O(d_D)$
coefficient operations; multiplication is a product of two coordinate vectors followed by reduction
modulo $F$, hence $O(d_D^2)$; inversion is an extended gcd in $\mathbb Q[t]/(F)$, also $O(d_D^2)$; each
coefficient operation is integer arithmetic on numbers of bit length $\poly(n,\tau)$. This is the standard
exact-arithmetic model of \cite{cohen93}, for which the input-size parameters of a field element are the
degree $d_D$, the height $h(F)$ of the defining polynomial and the height of the coordinates. All three
are $\poly(n,\tau)$ here, whereas the degree bound $d_D\le k^2$ alone controls none of the other two and
would by itself give no polynomial time bound.

\emph{Step 9 (kernels, subspaces and the Hadamard bound).} If $B$ is an $r\times r$ matrix with entries of
height at most $H$, Hadamard's inequality $|\det B|\le\prod_i\|b_i\|_2\le(\sqrt r\,2^{H})^{r}$ gives the
bound
\begin{equation}\label{eq:hadamard}
h(\det B)\ \le\ r\,H+\tfrac{r}{2}\log_2 r+O(r)
\end{equation}
used throughout. A basis of a kernel $D=\ker H$ or of an orthogonal complement is obtained by Cramer's
rule, so each of its entries is a quotient of two minors of order $r\le k$; the denominator is a non-zero
integer, hence of absolute value at least $1$, and by \eqref{eq:hadamard} the entry has height at most
$2\bigl(kH+\tfrac{k}{2}\log_2 k+O(k)\bigr)$. Consequently the bases of $D$, $D^{\perp_{K_0}}$, $T_D$ and
$\ker\bar\varphi$ inherit a factor $O(k)$ from the height of the matrix they come from, the functionals
$\varphi$, the prescribed image $\Phi_0$ and the entries of the cleared pencil are obtained from these by
polynomially many field operations, and their heights are again $\poly(n,\tau)$; clearing the chart
denominator $\delta$ multiplies them by at most $e\,h(\delta)=\poly(n,\tau)$, where $e\le\omega$ is the
clearing exponent.

\emph{Step 10 (the subresultant/gcd chain).} On a chart the maximal minors have $\lambda$-degree at most
$\omega\le k$, parameter degree at most $j\le2p$ and coefficient height at most $H=\poly(n,\tau)$
(Step~9). Since $p$ is fixed, a coefficient is a polynomial in the $p^2$ chart parameters with
$O((j+1)^{p^2})$ coefficients, a constant number, and we charge one ring operation on such a polynomial as
a single operation on integers of bit length $\varrho$. The gcd over $K(s)[\lambda]$ is computed by the
subresultant algorithm \cite{vzgg13}: two $\lambda$-polynomials of degree at most $\omega$ are reduced in
$O(\omega^2)$ ring operations, and every subresultant is a determinant of a Sylvester-type matrix of order
at most $2\omega$ whose entries are coefficients of the two inputs, so by \eqref{eq:hadamard} its height is
at most
\[
\varrho:=2\omega\,(H+\log_2\omega)+O(\omega),
\]
which is $\poly(n,\tau)$. The gcd divides a subresultant, so by Mignotte's bound (Step~6) its height exceeds that of the
subresultant by at most $O(\omega)$.
One gcd step therefore costs $O(\omega^2)$ ring operations on integers of bit length $\varrho$, i.e.\
$O(\omega^2\varrho^2)$ bit operations, a ring operation in $L_D$ carrying the extra factor $O(d_D^2)$ of
Step~$8'$. With the $\nu$ maximal minors of Claim~\ref{clm:degrees} the whole chain costs
\[
O\bigl(\nu\,\omega^2\,\varrho^2\,d_D^2\bigr)
\]
bit operations, and with $p$ fixed this is $\poly(n,\tau)$ because $\nu\le O(k^{2p+2})$.

Combining Steps 0--10 (with the inserted Steps $4'$ and $8'$) yields Claim~\ref{clm:arith} and the bound of
Theorem~\ref{thm:complexity}. No constant above is optimized and none is claimed to be sharp: each step
multiplies the height by a factor polynomial in $k$ and in the field degree $d_D\le k^2$, and the
additive term $\tfrac{r}{2}\log_2 r$ of \eqref{eq:hadamard} is kept as it stands. The two
quantities that we do not attempt to pin down are the multiplicative constant of the subresultant chain of
Step~10 and the exponent hidden in the factor $\poly(n,\tau)$ of the enumeration of
Claim~\ref{clm:arith}; both reach the final bound only through the coarse product displayed in the proof of
Theorem~\ref{thm:complexity}.

\normalsize

\begin{thebibliography}{99}
\bibitem{companion1} Deciding functional observer feasibility: compression, certificates, and
parameterized algorithms. Companion manuscript (preprint), 2026.
\bibitem{companion2} Minimum order and minimum-order construction for functional observers. Companion
manuscript (preprint), 2026.
\bibitem{companion3} Functional-observer feasibility in the hyperplane layer. Companion manuscript
(preprint), 2026.
\bibitem{fer26a} A. Fernando. \emph{Functional observer design: existence and order spectrum}.
arXiv:2609.33615, 2026 (preprint).
\bibitem{fer26b} A. Fernando. \emph{Functional observers: a chartwise characterisation}. arXiv:2609.37992,
2026 (preprint).
\bibitem{bruns88} W. Bruns, U. Vetter. \emph{Determinantal Rings}. Lecture Notes in Mathematics 1327.
Berlin, Germany: Springer, 1988.
\bibitem{fulton97} W. Fulton. \emph{Young Tableaux}. Cambridge, U.K.: Cambridge University Press, 1997
(degeneracy loci, Ch.~14).
\bibitem{bpr06} S. Basu, R. Pollack, M.-F. Roy. \emph{Algorithms in Real Algebraic Geometry}, 2nd ed.
Berlin, Germany: Springer, 2006.
\bibitem{cohen93} H. Cohen. \emph{A Course in Computational Algebraic Number Theory}. Graduate Texts in
Mathematics 138. Berlin, Germany: Springer, 1993 (exact arithmetic in number fields, resultants).
\bibitem{lll82} A. K. Lenstra, H. W. Lenstra, Jr., L. Lovász. Factoring polynomials with rational
coefficients. \emph{Math. Ann.}, vol.~261, no.~4, pp.~515--534, 1982.
\bibitem{vzgg13} J. von zur Gathen, J. Gerhard. \emph{Modern Computer Algebra}, 3rd ed. Cambridge, U.K.:
Cambridge University Press, 2013 (subresultants, gcd degree bounds).
\bibitem{emre80} E. Emre, M. L. J. Hautus. A polynomial characterization of $(A,B)$-invariant and
reachability subspaces. \emph{SIAM J. Control Optim.}, vol.~18, no.~4, 1980.
\bibitem{am69} M. F. Atiyah, I. G. Macdonald. \emph{Introduction to Commutative Algebra}.
Reading, MA, USA: Addison-Wesley, 1969 (integral dependence, integral closure, Ch.~5).
\end{thebibliography}

\end{document}
